\chapter{The Postulates of Quantum Mechanics}\label{ch:postulates}

The Introduction stated four rules of quantum mechanics informally, as a map of
what was to come. Chapter~\ref{ch:tools} built the language they are written
in, and this chapter states them precisely. The rules are called
\emph{postulates}: a postulate is not proved, it is the starting point from
which everything else is derived, and its justification is that its
consequences agree with what a laboratory records. The four postulates answer
four questions.

\begin{center}
\small
\begin{tabular}{@{}ll>{\raggedright\arraybackslash}p{0.35\textwidth}@{}}
  \toprule
  Postulate & Question it answers & Mathematical object\\
  \midrule
  1. State space & What is a state? & a unit vector in a Hilbert space $\Hil$\\
  2. Evolution & How does a closed system change? & a unitary $U$, or Schr\"odinger's equation\\
  3. Measurement & What can we learn, and at what cost? & measurement operators $\{M_m\}$\\
  4. Composition & How do systems combine? & the tensor product $\Hil_1\tp\Hil_2$\\
  \bottomrule
\end{tabular}
\end{center}

We state each postulate, work out what it means for the simplest quantum
system, the qubit, and then draw its consequences: the Bloch sphere,
quantum gates, the impossibility of telling non-orthogonal states apart, the
uncertainty principle, and entanglement. Three tools are introduced along the way because the postulates need them: positive operators (the operator version of a non-negative number), unitary representations of groups (the language of symmetries), and, building on the Riesz theorem of Section~\ref{sec:dual}, the description of a state as a functional on observables. The chapter ends where the subject
becomes an experimental science. We use entanglement to send two bits in one
qubit, extend the formalism to mixed states with density operators, and finish
with the argument of Einstein, Podolsky and Rosen and Bell's answer to it,
which shows that no classical story can reproduce what quantum mechanics predicts.

\begin{objectives}[In this chapter]
  You will state the four postulates and apply them to qubits; use the Bloch sphere and the
  standard gates; compute measurement probabilities and post-measurement states for
  projective, general and POVM measurements; recognise entangled states and work with
  tensor products; describe mixed states and subsystems with density operators;
  and derive the CHSH form of Bell's inequality and see why quantum mechanics violates it.
\end{objectives}

\begin{whereitsits}
  Every postulate uses a tool from Chapter~\ref{ch:tools}: states are unit vectors
  (Section~\ref{sec:hilbert}), evolution is unitary and measurement is projective or
  positive (Section~\ref{sec:operators}), and systems combine by the tensor product
  (Section~\ref{sec:tensor}). Chapter~\ref{ch:phenomena} applies the postulates to physical
  phenomena, and the rest of the programme (circuits, algorithms, error correction) is
  written in their language.
\end{whereitsits}

%======================================================================
\section{State space and evolution}\label{sec:state-evolution}
%======================================================================

\begin{objectives}
  You will state Postulates~1 and~2; describe a qubit state and place it on the Bloch sphere; use the Riesz representation theorem to see why every bra comes from a ket; recognise the basic single-qubit gates; pass between the unitary and the Schr\"odinger forms of the dynamics;
  and describe symmetries of a quantum system as unitary representations of groups.
\end{objectives}

\begin{recall}
  A Hilbert space, unit vectors and orthonormal bases (Section~\ref{sec:hilbert}); the Pauli matrices (Section~\ref{sec:vector-spaces}); Hermitian and unitary operators, the spectral theorem and the functional calculus $f(A)$ (Section~\ref{sec:operators}).
\end{recall}

Quantum mechanics is a \emph{framework}, not a theory. It does not tell us which Hilbert space describes a given physical system, or which Hamiltonian governs it; specific theories, such as quantum electrodynamics, answer those questions inside the framework. What the postulates fix is the \emph{form} of every such theory.
Do not expect them to be motivated: they emerged from decades of trial and error. Expect to learn how and when to apply them.

\subsection{State space}

\begin{postulate}[State space]\label{post:state}
  To every isolated physical system we associate a Hilbert space $\Hil$, its \term{state space}. The system is completely described by its \term{state vector}: a unit vector $\ket\psi\in\Hil$, that is, $\norm{\psi}=\sqrt{\braket\psi\psi}=1$.
\end{postulate}

Note what the postulate does \emph{not} say: it does not tell us which $\Hil$ describes a given system, nor which $\ket\psi$ the system is in. We simply posit a finite-dimensional state space $\C^d$ and work from there. Our fundamental system is the simplest nontrivial one, with $\Hil=\C^2$.

\begin{definition}[Qubit]\label{def:qubit}
  Fix an orthonormal basis $\{\ket0,\ket1\}$ of $\C^2$, the \term{computational basis}. A \term{qubit state} is
  \[
    \ket\psi=a\ket0+b\ket1,\qquad a,b\in\C,\qquad \braket\psi\psi=\abs a^2+\abs b^2=1 .
  \]
  The numbers $a,b$ are the \term{amplitudes} of $\ket0$ and $\ket1$, and $\ket\psi$ is a \term{superposition} of them.
\end{definition}

The vectors $\ket0,\ket1$ play the role of the two values of a classical bit. The difference is that superpositions $a\ket0+b\ket1$ also exist, and then it is meaningless to say that the qubit is ``really'' in $\ket0$ or in $\ket1$.
The unit-norm condition is what will make probabilities add up to $1$ (Section~\ref{sec:measurement}). For instance $\ket-=\frac1{\sqrt2}(\ket0-\ket1)$ has amplitude $1/\sqrt2$ for $\ket0$ and $-1/\sqrt2$ for $\ket1$.

\begin{example}\label{ex:normalise}
  For which $c>0$ is $\ket\psi=c\bigl(\ket0+(1-\ii)\ket1\bigr)$ a qubit state?
\end{example}
\begin{solution}
  \begin{steps}
    \item \emph{Write the norm.} $\braket\psi\psi=c^2\bigl(\abs1^2+\abs{1-\ii}^2\bigr)$.
    \item \emph{Evaluate.} $\abs{1-\ii}^2=1+1=2$, so $\braket\psi\psi=3c^2$.
    \item \emph{Solve.} $3c^2=1$ gives $c=1/\sqrt3$.
  \end{steps}
  Any multiple $c\,e^{\ii\gamma}$ with $\gamma\in\R$ also works. These describe the same physical state, as we will see in Section~\ref{sec:measurement}.
\end{solution}

\begin{example}[Reading the amplitudes]\label{ex:read-amplitudes}
  Is $\ket\psi=\tfrac15\bigl(3\ket0+4\ii\ket1\bigr)$ a state? What do its amplitudes say?
\end{example}
\begin{solution}
  \begin{steps}
    \item \emph{Normalisation.} $\abs{3/5}^2+\abs{4\ii/5}^2=\tfrac{9}{25}+\tfrac{16}{25}=1$: yes, it is a state.
    \item \emph{Moduli.} The squared moduli $\tfrac9{25}=0.36$ and $\tfrac{16}{25}=0.64$ will turn out to be the probabilities of finding $\ket0$ and $\ket1$ (Example~\ref{ex:meas-comp}).
    \item \emph{Phases.} The amplitude of $\ket1$ is $\tfrac45e^{\ii\pi/2}$: relative to the amplitude of $\ket0$ it carries a phase $\pi/2$. That relative phase does not change the probabilities in the computational basis, but it does change what other measurements will see.
  \end{steps}
\end{solution}

Count the parameters. The two amplitudes are four real numbers; normalisation removes one; and we will see that a global phase removes another. Only \emph{two} real parameters remain. Where do they live?

\subsection{States and linear functionals: the Riesz representation theorem}

In Section~\ref{sec:hilbert} we introduced the bra $\bra v$ as the row vector $\ket v^\dagger$, and observed that $\ket w\mapsto\braket vw$ is a linear map $\Hil\to\C$, that is, a \emph{linear functional} on $\Hil$. It is natural to ask whether we have found \emph{all} linear functionals this way. The answer is yes, and it is a theorem of Riesz.

\begin{theorem}[Riesz representation theorem, finite dimension]\label{thm:riesz}
  Let $\Hil$ be a finite-dimensional Hilbert space and $\varphi:\Hil\to\C$ a linear functional. Then there is a \emph{unique} vector $\ket v\in\Hil$ such that
  \[
    \varphi(\ket w)=\braket vw\qquad\text{for all }\ket w\in\Hil .
  \]
  Moreover $\abs{\varphi(\ket w)}\le\norm v\norm w$, and equality holds for $\ket w\propto\ket v$, so $\norm v$ is the ``size'' of $\varphi$.
\end{theorem}
\begin{proof}
  \begin{steps}
    \item \emph{Choose a basis.} Let $\{\ket{e_i}\}$ be an orthonormal basis of $\Hil$. Every $\ket w$ is $\sum_i\braket{e_i}w\ket{e_i}$ (Theorem~\ref{thm:parseval}), so by linearity $\varphi(\ket w)=\sum_i\braket{e_i}w\,\varphi(\ket{e_i})$.
    \item \emph{Define the candidate.} Put $\ket v=\sum_i\overline{\varphi(\ket{e_i})}\ket{e_i}$. Its bra is $\bra v=\sum_i\varphi(\ket{e_i})\bra{e_i}$ (conjugating the coefficients).
    \item \emph{Check.} $\braket vw=\sum_i\varphi(\ket{e_i})\braket{e_i}w=\varphi(\ket w)$ by the first step.
    \item \emph{Uniqueness.} If $\ket v,\ket{v'}$ both work, then $\braket{v-v'}w=0$ for all $\ket w$; taking $\ket w=\ket{v-v'}$ gives $\norm{v-v'}^2=0$.
    \item \emph{The bound.} It is Cauchy--Schwarz (Theorem~\ref{thm:cs}), with equality for $\ket w\propto\ket v$. \qedhere
  \end{steps}
\end{proof}

\begin{corollary}\label{cor:bras-kets}
  The map $\ket v\mapsto\bra v$ is a bijection from $\Hil$ onto the space $\Hil^*$ of linear functionals; it is conjugate-linear and preserves norms. So \emph{every bra is the bra of a ket}, which is what justifies treating bras and kets as two views of the same information.
\end{corollary}

\begin{example}[A functional on $\C^2$]\label{ex:riesz-c2}
  Find the vector representing $\varphi(z_1,z_2)=(1+\ii)z_1-2z_2$.
\end{example}
\begin{solution}
  \begin{steps}
    \item \emph{Values on the basis.} $\varphi(\ket0)=\varphi(1,0)=1+\ii$ and $\varphi(\ket1)=\varphi(0,1)=-2$.
    \item \emph{Representer.} $\ket v=\overline{(1+\ii)}\ket0+\overline{(-2)}\ket1=(1-\ii)\ket0-2\ket1$.
    \item \emph{Check.} $\braket vw=\overline{(1-\ii)}w_1+\overline{(-2)}w_2=(1+\ii)w_1-2w_2=\varphi(w)$.
  \end{steps}
\end{solution}

\begin{example}[Evaluation at a point, as an inner product]\label{ex:riesz-eval}
  On $P_1$ (polynomials of degree at most $1$) with $\langle p,q\rangle=\int_{-1}^1\overline p\,q\dd x$, represent the functional $\varphi(p)=p(0)$.
\end{example}
\begin{solution}
  \begin{steps}
    \item \emph{An orthonormal basis.} $e_0=\frac1{\sqrt2}$ and $e_1=\sqrt{\tfrac32}\,x$: indeed $\int_{-1}^1\tfrac12\dd x=1$, $\int_{-1}^1\tfrac32x^2\dd x=1$, $\int_{-1}^1x\dd x=0$.
    \item \emph{Values.} $\varphi(e_0)=\tfrac1{\sqrt2}$ and $\varphi(e_1)=0$.
    \item \emph{Representer.} $v=\overline{\varphi(e_0)}e_0+\overline{\varphi(e_1)}e_1=\tfrac1{\sqrt2}\cdot\tfrac1{\sqrt2}=\tfrac12$.
    \item \emph{Reading.} $p(0)=\int_{-1}^1\tfrac12\,p(x)\dd x$ for every $p\in P_1$: ``the value at the origin of a straight line is its average over $[-1,1]$''. \emph{Check} with $p=2+3x$: $p(0)=2$ and $\int_{-1}^1\tfrac12(2+3x)\dd x=2$.
  \end{steps}
\end{solution}

\begin{lookahead}
  Riesz's theorem has two uses in what follows. First, it is the reason a measurement outcome may be described by a ket: the functional ``amplitude to find $m$'' is the bra $\bra m$. Second, applied not to $\Hil$ but to the space of \emph{operators} with the Hilbert--Schmidt inner product, it will identify every linear functional on observables, and in particular every ``state'', with a density operator (Theorem~\ref{thm:riesz-operators}, Section~\ref{sec:density}). In infinite dimension the theorem still holds for \emph{continuous} functionals, and it is a pillar of the theory of Hilbert spaces.
\end{lookahead}

\subsection{The Bloch sphere}

Every qubit state, up to a global phase, can be written as
\begin{equation}\label{eq:bloch-param}
  \ket\psi=\cos\tfrac\theta2\,\ket0+e^{\ii\phi}\sin\tfrac\theta2\,\ket1,\qquad \theta\in[0,\pi],\ \phi\in[0,2\pi).
\end{equation}
To see this:
\begin{steps}
  \item multiply $\ket\psi=a\ket0+b\ket1$ by a phase so that the amplitude of $\ket0$ is real and non-negative;
  \item then $a=\cos\frac\theta2$ for some $\frac\theta2\in[0,\frac\pi2]$, and $\abs b=\sin\frac\theta2$ by normalisation;
  \item write $b=e^{\ii\phi}\abs b$.
\end{steps}
The two angles $\theta,\phi$ are the polar and azimuthal angles of a point of the unit sphere, so we associate with $\ket\psi$ its \term{Bloch vector}
\[
  \vec n=(\sin\theta\cos\phi,\ \sin\theta\sin\phi,\ \cos\theta)\in\R^3,\qquad \abs{\vec n}=1 .
\]

\begin{figure}[htb]
  \centering
  \begin{tikzpicture}[>=stealth,scale=1.25]
    % sphere outline and equator
    \draw[ink,thick] (0,0) circle (2);
    \draw[ink!60] (-2,0) arc (180:360:2 and 0.7);
    \draw[ink!40,dashed] (-2,0) arc (180:0:2 and 0.7);
    % axes
    \draw[->,ink!70] (0,0) -- (0,2.6) node[above] {$z$};
    \draw[->,ink!70] (0,0) -- (2.6,0) node[right] {$y$};
    \draw[->,ink!70] (0,0) -- (-1.1,-0.77) node[below left] {$x$};
    % named states
    \fill[navy] (0,2) circle (2pt) node[above left,font=\small] {$\ket0$};
    \fill[navy] (0,-2) circle (2pt) node[below right,font=\small] {$\ket1$};
    \fill[navy] (-0.9,-0.63) circle (2pt) node[below left,font=\small] {$\ket+$};
    \fill[navy] (0.9,0.63) circle (2pt) node[above right,font=\small] {$\ket-$};
    \fill[navy] (2,0) circle (2pt) node[above right,font=\small] {$\ket{+\ii}$};
    \fill[navy] (-2,0) circle (2pt) node[above left,font=\small] {$\ket{-\ii}$};
    % a generic state
    \draw[very thick,figacc,->] (0,0) -- (0.52,1.45);
    \fill[figacc] (0.52,1.45) circle (2pt) node[right,font=\small] {$\ket\psi$};
    \draw[figacc] (0,0.55) arc (90:70:0.55) node[midway,above,font=\scriptsize] {$\theta$};
  \end{tikzpicture}
  \caption{The Bloch sphere. The six states $\ket0,\ket1$ (poles), $\ket\pm$ (along the $x$-axis) and $\ket{\pm\ii}$ (along the $y$-axis) are marked.
  A general state $\ket\psi$ is a point of the sphere; orthogonal states sit at antipodal points.}
  \label{fig:bloch}
\end{figure}

The poles are $\ket0$ and $\ket1$; $\ket\pm$ lie on the $\pm x$ axis and $\ket{\pm\ii}$ on the $\pm y$ axis (Figure~\ref{fig:bloch}). The Bloch vector can be read off from the matrix $\ket\psi\!\bra\psi$:

\begin{proposition}\label{prop:bloch}
  For $\ket\psi$ as in \eqref{eq:bloch-param} with Bloch vector $\vec n$,
  \[
    \ket\psi\!\bra\psi=\tfrac12\bigl(\Id+\vec n\cdot\vec\sigma\bigr),\qquad n_j=\bra\psi\sigma_j\ket\psi .
  \]
  Moreover, if $\ket\varphi$ has Bloch vector $\vec m$, then $\abs{\braket\varphi\psi}^2=\tfrac12(1+\vec m\cdot\vec n)$.
\end{proposition}
\begin{proof}
  Write $c=\cos\frac\theta2$, $s=\sin\frac\theta2$.
  \begin{steps}
    \item \emph{The projector.} $\ket\psi\!\bra\psi=\begin{pmatrix}c^2&cs\,e^{-\ii\phi}\\cs\,e^{\ii\phi}&s^2\end{pmatrix}=\frac12\begin{pmatrix}1+\cos\theta&\sin\theta\,e^{-\ii\phi}\\\sin\theta\,e^{\ii\phi}&1-\cos\theta\end{pmatrix}$, using $c^2=\tfrac12(1+\cos\theta)$, $s^2=\tfrac12(1-\cos\theta)$, $cs=\tfrac12\sin\theta$.
    \item \emph{The Pauli side.} $\vec n\cdot\vec\sigma=\begin{pmatrix}n_z&n_x-\ii n_y\\n_x+\ii n_y&-n_z\end{pmatrix}$, so $\tfrac12(\Id+\vec n\cdot\vec\sigma)$ has the same four entries. This proves the first identity.
    \item \emph{The components.} By $\Tr(\sigma_j\sigma_k)=2\delta_{jk}$ and $\Tr\sigma_j=0$: $\bra\psi\sigma_j\ket\psi=\Tr(\sigma_j\ket\psi\!\bra\psi)=n_j$.
    \item \emph{The overlap.} $\abs{\braket\varphi\psi}^2=\Tr\bigl(\ket\varphi\!\bra\varphi\ket\psi\!\bra\psi\bigr)=\tfrac14\Tr\bigl((\Id+\vec m\cdot\vec\sigma)(\Id+\vec n\cdot\vec\sigma)\bigr)=\tfrac14(2+2\vec m\cdot\vec n)$. \qedhere
  \end{steps}
\end{proof}

In particular, $\vec m=-\vec n$ gives overlap $0$: \emph{orthogonal states are antipodal points of the sphere.} So a qubit is a point on a sphere; gates will turn out to be rotations of it.

\begin{example}\label{ex:bloch-vectors}
  Find the Bloch vectors of $\ket0,\ket1,\ket\pm,\ket{\pm\ii}$, and of $\ket\psi=\frac1{\sqrt3}\ket0+\sqrt{\tfrac23}\,e^{\ii\pi/3}\ket1$.
\end{example}
\begin{solution}
  \begin{steps}
    \item \emph{The basis states.} Use $n_j=\bra\psi\sigma_j\ket\psi$. $\ket0$ gives $(0,0,1)$ and $\ket1$ gives $(0,0,-1)$.
    \item \emph{The $x$ states.} $\sx\ket\pm=\pm\ket\pm$, and $\bra\pm\sz\ket\pm=\bra\pm\sy\ket\pm=0$, so $\ket\pm\mapsto(\pm1,0,0)$. Likewise $\ket{\pm\ii}\mapsto(0,\pm1,0)$.
    \item \emph{The general state: match \eqref{eq:bloch-param}.} $\cos\frac\theta2=\frac1{\sqrt3}$, so $\cos\theta=2\cos^2\frac\theta2-1=-\tfrac13$ and $\sin\theta=\tfrac{2\sqrt2}{3}$; and $\phi=\pi/3$.
    \item \emph{Assemble.} $\vec n=\bigl(\tfrac{2\sqrt2}{3}\cdot\tfrac12,\ \tfrac{2\sqrt2}{3}\cdot\tfrac{\sqrt3}{2},\ -\tfrac13\bigr)=\bigl(\tfrac{\sqrt2}{3},\ \tfrac{\sqrt6}{3},\ -\tfrac13\bigr)$.
    \item \emph{Check.} $\tfrac29+\tfrac69+\tfrac19=1$.
  \end{steps}
\end{solution}

\begin{example}[Overlap from the geometry]\label{ex:bloch-overlap}
  Compute $\abs{\braket0{+}}^2$, $\abs{\braket{+}{+\ii}}^2$ and $\abs{\braket{+}{-}}^2$ using Bloch vectors.
\end{example}
\begin{solution}
  \begin{steps}
    \item \emph{$\ket0,\ket+$:} $\vec m\cdot\vec n=(0,0,1)\cdot(1,0,0)=0$, so the overlap is $\tfrac12(1+0)=\tfrac12$. (Directly: $\braket0+=\tfrac1{\sqrt2}$.)
    \item \emph{$\ket+,\ket{+\ii}$:} $(1,0,0)\cdot(0,1,0)=0$, overlap $\tfrac12$.
    \item \emph{$\ket+,\ket-$:} $(1,0,0)\cdot(-1,0,0)=-1$, overlap $\tfrac12(1-1)=0$: orthogonal, antipodal.
  \end{steps}
  In the picture of Figure~\ref{fig:bloch}, the squared overlap is determined by the angle between the two vectors alone.
\end{solution}

\subsection{Unitary evolution}

\begin{postulate}[Evolution]\label{post:evolution}
  The evolution of a \emph{closed} quantum system is described by a unitary transformation. The state $\ket\psi$ at time $t_1$ and the state $\ket{\psi'}$ at time $t_2$ are related by a unitary operator $U$ that depends only on $t_1$ and $t_2$:
  \[
    \ket{\psi'}=U\ket\psi .
  \]
\end{postulate}

\emph{Closed} means not interacting with any other system. Strictly speaking, only the whole Universe is closed, but many systems are closed to excellent approximation, and in principle every open system is part of a larger closed one.
Again the postulate only guarantees that the evolution can be written this way: it does not say which $U$ describes a real system. Unitarity is exactly what is needed to keep Postulate~1 consistent: $\braket{\psi'}{\psi'}=\bra\psi U^\dagger U\ket\psi=1$, so states are sent to states.
For a single qubit, any unitary can be realised in realistic systems, and the standard ones have names.

\begin{center}
\small
\begin{tabular}{@{}lcll@{}}
  \toprule
  Gate & Matrix & Action & Name\\
  \midrule
  $X=\sx$ & $\begin{pmatrix}0&1\\1&0\end{pmatrix}$ & $\ket0\leftrightarrow\ket1$ & bit flip, quantum NOT\\[6pt]
  $Z=\sz$ & $\begin{pmatrix}1&0\\0&-1\end{pmatrix}$ & $\ket0\mapsto\ket0,\ \ket1\mapsto-\ket1$ & phase flip\\[6pt]
  $H$ & $\frac1{\sqrt2}\begin{pmatrix}1&1\\1&-1\end{pmatrix}$ & $\ket0\mapsto\ket+,\ \ket1\mapsto\ket-$ & Hadamard\\[6pt]
  $S$ & $\begin{pmatrix}1&0\\0&\ii\end{pmatrix}$ & $\ket0\mapsto\ket0,\ \ket1\mapsto\ii\ket1$ & phase gate\\
  \bottomrule
\end{tabular}
\end{center}

$X$ and $Z$ are Hermitian and unitary (they are Pauli matrices, Proposition~\ref{prop:pauli}). The Hadamard gate turns ``definite'' states into equal superpositions, which is the first step of most quantum algorithms. $S$ is unitary but not Hermitian.

\begin{pitfall}
  The symbol $H$ stands for two different things: the Hadamard gate above, and the Hamiltonian of the next subsection. Context decides which.
\end{pitfall}

\begin{example}\label{ex:hadamard-eig}
  Verify that the Hadamard gate $H$ is unitary with $H^2=\Id$, and find its eigenvalues and eigenvectors.
\end{example}
\begin{solution}
  \begin{steps}
    \item \emph{Unitarity.} By Example~\ref{ex:hadamard}, $H=(\sx+\sz)/\sqrt2$ is Hermitian with $H^2=\Id$, so $H^\dagger H=H^2=\Id$.
    \item \emph{Eigenvalues.} Write $H=\vec n\cdot\vec\sigma$ with the unit vector $\vec n=\frac1{\sqrt2}(1,0,1)$. Then $H^2=\Id$ (product formula of Proposition~\ref{prop:pauli}), so the eigenvalues are $\pm1$ (Exercise~\ref{ex:vn-sigma-eig}).
    \item \emph{Eigenvector for $+1$.} Solve $(H-\Id)v=0$: the first row gives $v_1(1-\sqrt2)+v_2=0$, so $v\propto(1,\sqrt2-1)=(1,\tan\frac\pi8)\propto\bigl(\cos\tfrac\pi8,\sin\tfrac\pi8\bigr)$.
    \item \emph{Eigenvector for $-1$.} It must be orthogonal to the first: $\bigl(-\sin\tfrac\pi8,\cos\tfrac\pi8\bigr)$.
    \item \emph{Geometry.} These are the Bloch vectors $\pm\vec n$, at angle $\pi/4$ from $\ket0$: the Hadamard gate is a half-turn of the sphere about the axis $\vec n$.
  \end{steps}
\end{solution}

\begin{example}[A circuit of gates]\label{ex:circuit}
  Start in $\ket0$, apply $H$, then $S$, then $H$. Compare the final state with the one obtained without the middle gate.
\end{example}
\begin{solution}
  \begin{steps}
    \item \emph{After $H$:} $\ket+$.
    \item \emph{After $S$:} $S\ket+=\frac1{\sqrt2}(\ket0+\ii\ket1)=\ket{+\ii}$.
    \item \emph{After the last $H$:} $H\ket{+\ii}=\tfrac1{\sqrt2}\bigl(H\ket0+\ii H\ket1\bigr)=\tfrac1{\sqrt2}\bigl(\ket++\ii\ket-\bigr)=\tfrac12\bigl((1+\ii)\ket0+(1-\ii)\ket1\bigr)$.
    \item \emph{Without $S$:} $HH\ket0=\ket0$.
  \end{steps}
  The final amplitudes have squared moduli $\tfrac{\abs{1\pm\ii}^2}{4}=\tfrac12$ each, instead of $(1,0)$. The phase gate $S$ is invisible if one looks at $\ket0,\ket1$ probabilities straight after it, but it changes the outcome after a second Hadamard: relative phases are observable (Section~\ref{sec:measurement}).
\end{solution}

\subsection{The Schr\"odinger equation}

In continuous time, the dynamics is governed by an operator that stores the energy.

\begin{proposition}[Evolution, Schr\"odinger form]\label{prop:schrodinger}
  Postulate~\ref{post:evolution} is equivalent to the following statement. The state of a closed system obeys
  \begin{equation}\label{eq:schrodinger-abstract}
    \ii\hbar\,\frac{d}{dt}\ket{\psi(t)}=H\ket{\psi(t)},
  \end{equation}
  where $\hbar$ is Planck's constant (fixed experimentally) and $H=H^\dagger$ is a fixed Hermitian operator, the \term{Hamiltonian} of the system. The solution is $\ket{\psi(t_2)}=U(t_1,t_2)\ket{\psi(t_1)}$ with
  \[
    U(t_1,t_2)=\exp\!\Bigl(-\frac{\ii H(t_2-t_1)}{\hbar}\Bigr).
  \]
\end{proposition}
\begin{proof}
  \begin{steps}
    \item \emph{From $H$ to $U$.} If $H$ is Hermitian, so is $K=-H(t_2-t_1)/\hbar$, and $U(t_1,t_2)=e^{\ii K}$ is unitary by Proposition~\ref{prop:expunitary}.
    \item \emph{It solves the equation.} $\frac{d}{dt_2}U=-\frac{\ii H}{\hbar}U$, so $U\ket{\psi(t_1)}$ satisfies \eqref{eq:schrodinger-abstract}.
    \item \emph{From $U$ to $H$.} Let $U$ be unitary. By Theorem~\ref{thm:spectrum}(b) and the spectral theorem, $U=\sum_je^{\ii\theta_j}\ket{v_j}\!\bra{v_j}$ with $\theta_j\in\R$.
    \item \emph{Take a logarithm.} Put $K=\sum_j\theta_j\ket{v_j}\!\bra{v_j}$, Hermitian, so that $U=e^{\ii K}$; then $U=e^{-\ii H\Delta t/\hbar}$ for the Hermitian operator $H=-\hbar K/\Delta t$. \qedhere
  \end{steps}
\end{proof}

The correspondence between unitaries and Hamiltonians is one-to-one up to the choice of branch of the logarithm ($\theta_j$ is determined only modulo $2\pi$). Knowing $H$ and $\hbar$ we know the dynamics completely, in principle; finding the right $H$ for a physical system is hard and needs experiments.
In quantum computing we usually just posit a matrix $H$, and often set $\hbar=1$.

Since $H$ is Hermitian, it has a spectral decomposition $H=\sum_EE\ket E\!\bra E$. The eigenvectors $\ket E$ are the \term{energy eigenstates} or stationary states, $E$ is the energy, the smallest $E$ is the \emph{ground state energy} and its eigenspace the ground state. By the functional calculus,
\[
  e^{-\ii Ht/\hbar}\ket E=e^{-\ii Et/\hbar}\ket E :
\]
a stationary state only picks up a global phase, so nothing observable changes. A \emph{superposition} of energy levels, on the other hand, develops \emph{relative} phases between the levels, and those are observable.

\begin{example}\label{ex:H-eq-X}
  A qubit has Hamiltonian $H=\hbar\omega X$ and starts in $\ket{\psi(0)}=\ket0$. Find the energy levels, the state at time $t$, and the probability of finding $\ket1$.
\end{example}
\begin{solution}
  \begin{steps}
    \item \emph{Energy levels.} $X\ket\pm=\pm\ket\pm$, so the energy eigenstates are $\ket\pm$ with energies $\pm\hbar\omega$; the ground state is $\ket-$.
    \item \emph{The propagator.} Using $X^2=\Id$ and Example~\ref{ex:funcs}(b), $U(t)=e^{-\ii\omega tX}=\cos(\omega t)\,\Id-\ii\sin(\omega t)\,X$.
    \item \emph{The state.} $\ket{\psi(t)}=U(t)\ket0=\cos(\omega t)\ket0-\ii\sin(\omega t)\ket1$.
    \item \emph{The amplitude of $\ket1$.} It is $-\ii\sin\omega t$; its modulus squared, which will be a probability by Postulate~3, is $\sin^2(\omega t)$.
    \item \emph{Interpretation.} The initial state $\ket0=\frac1{\sqrt2}(\ket++\ket-)$ is a superposition of the two energy levels, and the relative phase $e^{-2\ii\omega t}$ between them is what makes the weight on $\ket1$ oscillate, with period $\pi/\omega$. Had we started in $\ket+$, nothing would have happened.
  \end{steps}
\end{solution}

\begin{example}[A precessing qubit]\label{ex:precession}
  Let $H=\hbar\omega Z$ and $\ket{\psi(0)}=\ket+$. Find $\ket{\psi(t)}$ and its Bloch vector.
\end{example}
\begin{solution}
  \begin{steps}
    \item \emph{The energy eigenstates} are $\ket0,\ket1$ with energies $+\hbar\omega,-\hbar\omega$.
    \item \emph{Evolve each component.} $\ket{\psi(0)}=\frac1{\sqrt2}(\ket0+\ket1)$ goes to $\ket{\psi(t)}=\frac1{\sqrt2}\bigl(e^{-\ii\omega t}\ket0+e^{\ii\omega t}\ket1\bigr)$.
    \item \emph{Factor a global phase.} $\ket{\psi(t)}=e^{-\ii\omega t}\cdot\frac1{\sqrt2}\bigl(\ket0+e^{2\ii\omega t}\ket1\bigr)$.
    \item \emph{Read off $\theta,\phi$.} Compare with \eqref{eq:bloch-param}: $\theta=\pi/2$ and $\phi=2\omega t$, so $\vec n(t)=(\cos2\omega t,\ \sin2\omega t,\ 0)$.
  \end{steps}
  The Bloch vector rotates in the equator at angular velocity $2\omega$: the qubit \emph{precesses} about the $z$-axis, exactly as a classical magnet does about a field (the ``Larmor precession'').
\end{solution}

\begin{remark}
  A puzzle: in circuits we ``apply the gate $X$ to a qubit''. But applying means that an external agent interacts with the qubit, so it is not closed, which seems to contradict Postulate~2. The resolution is that the atom plus the laser form a closed system, and, looking at the atom alone, its state evolves almost exactly as if governed by a Hamiltonian $H(t)$ that depends on parameters the experimentalist controls
  (the laser's intensity and duration, for example). We will keep describing the evolution of quantum systems, even those that are not closed, by unitaries. The main exception is measurement, to which we turn next.
\end{remark}

\subsection{Symmetries and unitary representations of groups}\label{sec:grouprep}

Gates form a structure: doing $X$ twice is the identity, doing a rotation by $\alpha$ and then by $\beta$ is a rotation by $\alpha+\beta$, and every operation can be undone. This is the structure of a group, and the right way to organise it is the language of \emph{representations}.

A \term{group} is a set $G$ with an associative product, an identity $e$, and an inverse $g^{-1}$ for each element. A \term{symmetry} of a physical system is a transformation of the system (a rotation of the laboratory, a relabelling of identical particles, a shift in time) that does not change any measurable probability.
Wigner's theorem says that each symmetry acts on state vectors by a unitary operator (or an anti-unitary one, for time reversal), up to a global phase. Composing symmetries then composes the unitaries. This is the following concept.

\begin{definition}[Unitary representation]\label{def:unitaryrep}
  A \term{unitary representation} of a group $G$ on a Hilbert space $\Hil$ is a map $g\mapsto U(g)$ from $G$ into the unitary operators on $\Hil$ such that
  \[
    U(gh)=U(g)\,U(h)\quad\text{for all }g,h\in G,\qquad U(e)=\Id .
  \]
\end{definition}

Two consequences are immediate:
\begin{steps}
  \item $U(g^{-1})=U(g)^{-1}=U(g)^\dagger$: undoing a symmetry is the inverse of its unitary. (From $U(g)U(g^{-1})=U(e)=\Id$.)
  \item if $G$ is \emph{abelian}, all the $U(g)$ commute, and being unitary (hence normal) they have a common orthonormal eigenbasis (Exercise~\ref{ex:commuting-basis}). Every unitary representation of an abelian group splits the space into one-dimensional pieces on which $U(g)$ acts as a phase $\chi(g)$, with $\abs{\chi(g)}=1$.
\end{steps}

\begin{recall}
  Examples (e) below and Example~\ref{ex:symmetric-H} use the tensor product, which is treated properly in Section~\ref{sec:composite}. All we need for now: $\C^2\otimes\C^2=\C^4$ has the orthonormal basis $\ket a\ket b=\ket a\otimes\ket b$ with $a,b\in\{0,1\}$; $(A\otimes B)(\ket a\ket b)=A\ket a\otimes B\ket b$; and the matrix of $A\otimes B$ is the Kronecker product of the two matrices.
\end{recall}

\begin{example}[Examples of representations]\label{ex:reps}
  Check that each of the following is a unitary representation.
  \begin{enumerate}
    \item[(a)] $G=\{e,g\}\cong\Z_2$ on $\C^2$, with $U(e)=\Id$, $U(g)=X$.
    \item[(b)] The circle group $G=\{e^{\ii\varphi}\}$ (or $\R$ mod $2\pi$) on $\C^2$, with $U(\varphi)=e^{\ii\varphi Z}=\diag(e^{\ii\varphi},e^{-\ii\varphi})$.
    \item[(c)] $G=(\R,+)$ on any $\Hil$, with $U(t)=e^{-\ii Ht/\hbar}$ for a Hermitian $H$ (time translations).
    \item[(d)] $G=\Z_N$ on $\C^N$ with $U(k)=S^k$, where $S\ket j=\ket{j+1\bmod N}$ is the cyclic shift.
    \item[(e)] $G=S_2=\{e,\tau\}$ on $\C^2\otimes\C^2$ with $U(\tau)=\mathrm{SWAP}$, $\mathrm{SWAP}\ket a\ket b=\ket b\ket a$.
  \end{enumerate}
\end{example}
\begin{solution}
  In each case we check the three requirements: unitarity, $U(e)=\Id$, and the product rule.
  \begin{steps}
    \item (a) $X$ is unitary and $X^2=\Id=U(g^2)=U(e)$: the product rule holds for every pair ($g\cdot g=e$).
    \item (b) $e^{\ii\varphi Z}$ is unitary (Proposition~\ref{prop:expunitary}); $\diag(e^{\ii\varphi},e^{-\ii\varphi})\diag(e^{\ii\varphi'},e^{-\ii\varphi'})=\diag(e^{\ii(\varphi+\varphi')},e^{-\ii(\varphi+\varphi')})$, which is $U(\varphi+\varphi')$.
    \item (c) $U(t)$ is unitary since $H$ is Hermitian, and $U(t+s)=U(t)U(s)$ because $-\ii Ht/\hbar$ and $-\ii Hs/\hbar$ commute (Proposition~\ref{prop:commuting-exp}).
    \item (d) $S$ merely permutes the orthonormal basis, so it is unitary, and $S^N=\Id$ ensures that $U(k)=S^k$ depends only on $k\bmod N$. The product rule is $S^{k+l}=S^kS^l$.
    \item (e) $\mathrm{SWAP}$ permutes the orthonormal basis $\ket{ab}$, so it is unitary, and $\mathrm{SWAP}^2=\Id$.
  \end{steps}
\end{solution}

The representation (c) deserves a comment. The group $(\R,+)$ is \emph{time translation}, and the Schr\"odinger equation says that closed-system dynamics is a unitary representation of it. \term{Stone's theorem} states the converse: every (continuous) unitary representation of $(\R,+)$ has the form $e^{-\ii Ht/\hbar}$ for a unique Hermitian $H$. So \emph{a Hamiltonian is the same thing as a one-parameter group of unitaries}, which is how Proposition~\ref{prop:schrodinger} generalises to infinite dimension.

\paragraph{The qubit and the rotation group.}
The most important representation in this book is that of rotations. A rotation of the Bloch sphere by an angle $\theta$ about the axis $\vec n$ is implemented on states by the unitary $R_{\vec n}(\theta)=e^{-\ii\theta\vec n\cdot\vec\sigma/2}$ (Exercise~\ref{ex:rotation-formula}), and
$R_{\vec n}(\theta)R_{\vec n}(\theta')=R_{\vec n}(\theta+\theta')$. But, curiously,
\[
  R_{\vec n}(\theta+2\pi)=-R_{\vec n}(\theta),\qquad\text{so}\qquad R_{\vec n}(2\pi)=-\Id .
\]
A rotation by $2\pi$ is the identity on the Bloch sphere, but acts on the \emph{state vector} as $-1$: a global phase. This is perfectly legitimate, because global phases are unobservable (Section~\ref{sec:measurement}); so the qubit carries a representation of the rotation group \emph{up to sign}, that is,
a representation of its double cover $SU(2)$, the group of $2\times2$ unitary matrices of determinant $1$, acting on $\C^2$ by $U(g)=g$.

\begin{example}[Symmetries of a two-qubit Hamiltonian]\label{ex:symmetric-H}
  Let $H=\hbar J\,(X\otimes X+Y\otimes Y+Z\otimes Z)$ with $J>0$. Show that $H=\hbar J(2\,\mathrm{SWAP}-\Id)$, find its eigenvalues and eigenspaces, and show that $H$ commutes with the swap and with $U\otimes U$ for every $U\in SU(2)$.
\end{example}
\begin{solution}
  \begin{steps}
    \item \emph{The matrices.} In the basis $\ket{00},\ket{01},\ket{10},\ket{11}$: $Z\otimes Z=\diag(1,-1,-1,1)$; $X\otimes X$ has ones on the anti-diagonal; $Y\otimes Y$ has $-1$ at positions $(1,4),(4,1)$ and $+1$ at $(2,3),(3,2)$.
    \item \emph{Add them.} The corner entries cancel ($1-1=0$) and the middle block becomes $\begin{pmatrix}-1&2\\2&-1\end{pmatrix}$:
      \[
        X\otimes X+Y\otimes Y+Z\otimes Z=\begin{pmatrix}1&0&0&0\\0&-1&2&0\\0&2&-1&0\\0&0&0&1\end{pmatrix}=2\,\mathrm{SWAP}-\Id,
      \]
      since $\mathrm{SWAP}$ is the matrix with ones at $(1,1),(2,3),(3,2),(4,4)$.
    \item \emph{Eigenvalues.} $\mathrm{SWAP}$ has eigenvalue $+1$ on the three-dimensional \emph{symmetric} space spanned by $\ket{00},\ket{11},\frac1{\sqrt2}(\ket{01}+\ket{10})$, and $-1$ on the \emph{antisymmetric} vector $\ket{\Psi^-}=\frac1{\sqrt2}(\ket{01}-\ket{10})$. So $H$ has eigenvalue $\hbar J$ (three-fold) on the symmetric space and $-3\hbar J$ on the singlet. The ground state is the singlet.
    \item \emph{Symmetry.} $H$ is a combination of $\mathrm{SWAP}$ and $\Id$, so it commutes with $\mathrm{SWAP}$. It also commutes with $U\otimes U$: $\mathrm{SWAP}(U\otimes U)=(U\otimes U)\mathrm{SWAP}$ because both sides send $\ket a\ket b$ to $U\ket b\otimes U\ket a$.
  \end{steps}
  Notice what this says about the ground state: the singlet is invariant, up to $\det U=1$, under every $U\otimes U$ (equation~\eqref{eq:singlet-inv} below). It is the unique vector fixed by the whole group $SU(2)$ acting diagonally, which is the sense in which it has ``total spin zero''.
\end{solution}

\begin{proposition}[Symmetries preserve energy levels]\label{prop:sym-energy}
  If $U(g)H=HU(g)$ for all $g$ (a \emph{symmetry of the dynamics}), then each $U(g)$ maps every eigenspace of $H$ into itself.
\end{proposition}
\begin{proof}
  If $H\ket v=E\ket v$ then $H\,U(g)\ket v=U(g)H\ket v=E\,U(g)\ket v$. So $U(g)\ket v$ is again an eigenvector with the same energy. \qedhere
\end{proof}

\begin{lookahead}
  Symmetries organise quantum mechanics. In Example~\ref{ex:symmetric-H} the degeneracy of the symmetric levels is \emph{forced} by the symmetry group; in Section~\ref{sec:bell} the invariance of the singlet under $U\otimes U$ is what makes its correlations independent of the direction of measurement; and in quantum computing the Pauli group $\{\pm\Id,\pm\ii\Id\}\times\{\Id,X,Y,Z\}$ is the backbone of error correction. Finally, the discrete Fourier transform of Section~\ref{sec:fourier} is the change of basis that diagonalises the cyclic shift of Example~\ref{ex:reps}(d) (Exercise~\ref{ex:shift-fourier}): it is the decomposition of the representation of $\Z_N$ into its one-dimensional pieces.
\end{lookahead}

\subsection{Exponentials of non-commuting operators}

Gates are exponentials of Hamiltonians, so we must know when exponentials combine.

\begin{proposition}\label{prop:commuting-exp}
  If $A,B$ are commuting Hermitian operators, then $e^{\ii A}e^{\ii B}=e^{\ii(A+B)}$.
\end{proposition}
\begin{proof}
  \begin{steps}
    \item Commuting Hermitian operators can be diagonalised simultaneously: there is one orthonormal basis $\{\ket{v_i}\}$ with $A\ket{v_i}=a_i\ket{v_i}$ and $B\ket{v_i}=b_i\ket{v_i}$ (Exercise~\ref{ex:commuting-basis}).
    \item Both sides then act on $\ket{v_i}$ as multiplication by $e^{\ii(a_i+b_i)}$, so they agree on a basis. \qedhere
  \end{steps}
\end{proof}

If $A,B$ do not commute, the law fails: for $A=s\sx$, $B=s\sz$ and a small $s\neq0$, $e^{\ii A}e^{\ii B}\neq e^{\ii(A+B)}$ (compare second-order terms; Exercise~\ref{ex:trotter-exercise}). The general statement is the Baker--Campbell--Hausdorff formula,
$e^Ae^B=\exp\bigl(A+B+\tfrac12[A,B]+\cdots\bigr)$, where the dots are nested commutators. The failure is small when the operators are small, which gives the useful \emph{Trotter formula}
\[
  e^{\ii(A+B)}=\lim_{n\to\infty}\bigl(e^{\ii A/n}e^{\ii B/n}\bigr)^n .
\]
A Hamiltonian $H=H_1+H_2$ can therefore be simulated by alternating the evolutions under $H_1$ and under $H_2$, each of which may be easy: this is the idea behind quantum simulation on a quantum computer.

\begin{example}\label{ex:trotter}
  Let $A=\frac\pi4\sx$ and $B=\frac\pi4\sz$. Compare $e^{-\ii(A+B)}$ with $\bigl(e^{-\ii A/n}e^{-\ii B/n}\bigr)^n$.
\end{example}
\begin{solution}
  \begin{steps}
    \item \emph{The exact value.} $A+B=\frac\pi4(\sx+\sz)=\frac{\pi\sqrt2}{4}H$, where $H=\vec n\cdot\vec\sigma$ with $\vec n$ a unit vector and $H^2=\Id$.
    \item \emph{Exponentiate.} By the same computation as in Example~\ref{ex:funcs}(b), $e^{-\ii(A+B)}=\cos\tfrac{\pi}{2\sqrt2}\,\Id-\ii\sin\tfrac{\pi}{2\sqrt2}\,H$ (Exercise~\ref{ex:rotation-formula}).
    \item \emph{The Trotter products.} Computing numerically, the operator-norm error of $\bigl(e^{-\ii A/n}e^{-\ii B/n}\bigr)^n$ is about $0.54$ for $n=1$, $0.25$ for $n=2$, $0.050$ for $n=10$ and $0.0050$ for $n=100$.
    \item \emph{Conclusion.} The error decreases like $1/n$: ten times more steps buy ten times more accuracy.
  \end{steps}
\end{solution}

\begin{exerciseset}
  \exercise{1} Find $c>0$ such that $c\bigl((1+\ii)\ket0+2\ket1\bigr)$ is a qubit state. Which of these pairs of states differ only by a global phase: $\ket\psi$ and $-\ket\psi$; $\frac1{\sqrt2}(\ket0+\ket1)$ and $\frac1{\sqrt2}(\ket0+\ii\ket1)$; $\frac1{\sqrt2}(\ket0+\ii\ket1)$ and $\frac1{\sqrt2}(\ii\ket0-\ket1)$?

  \exercise{1} Which of $\frac1{\sqrt2}(\ket0+\ket1)$, $\frac12(\ket0+\ket1)$, $\frac{1}{\sqrt{13}}(2\ket0+3\ii\ket1)$ and $\frac13(2\ket0+\ii\ket1)$ are qubit states? For those that are, find the squared moduli of the amplitudes.

  \exercise{2} Prove the Riesz representation theorem in $\C^n$ with the standard inner product directly: if $\varphi(w)=\sum_ic_iw_i$ then $v=(\bar c_1,\dots,\bar c_n)$. Compute the representer of $\varphi(z_1,z_2,z_3)=z_1-\ii z_2+(2+\ii)z_3$ and its norm.

  \exercise{2} On $P_2$ with $\langle p,q\rangle=\int_{-1}^1\overline pq\dd x$, find the vector representing $\varphi(p)=p(0)$ and the one representing $\varphi(p)=p'(0)$. (Use the orthonormal basis $\frac1{\sqrt2},\sqrt{\tfrac32}x,\sqrt{\tfrac58}(3x^2-1)$.)

  \exercise{2} Show that the representer $v$ in Theorem~\ref{thm:riesz} does not depend on the orthonormal basis used to construct it. Show that $\ket v\mapsto\bra v$ is conjugate-linear and preserves norms (Corollary~\ref{cor:bras-kets}).

  \exercise{1} Compute the Bloch vectors of $\ket0,\ket1,\ket\pm,\ket{\pm\ii}$, and find $\theta,\phi$ and $\vec n$ for $\ket\psi=\frac1{\sqrt3}\ket0+\sqrt{\tfrac23}\,e^{\ii\pi/3}\ket1$.

  \exercise{2} Show that $\ket\psi\!\bra\psi=\frac12(\Id+\vec n\cdot\vec\sigma)$ and $\abs{\braket\varphi\psi}^2=\frac12(1+\vec m\cdot\vec n)$ (Proposition~\ref{prop:bloch}), and deduce that orthogonal states have antipodal Bloch vectors.

  \exercise{2} For $\ket\psi=a\ket0+b\ket1$ with $a\neq0$, find $\gamma$ such that $e^{-\ii\gamma}\ket\psi$ has a real, positive $\ket0$-amplitude, and read off $\theta/2$ and $\phi$. Show that two qubit states $\ket\psi,\ket\varphi$ describe the same Bloch vector if and only if $\abs{\braket\psi\varphi}=1$.

  \exercise{1} Verify that the Hadamard gate is unitary, that $H^2=\Id$, and find its eigenvalues and eigenvectors. (Hint: $H=\vec n\cdot\vec\sigma$ with $\vec n=\frac1{\sqrt2}(1,0,1)$.)\label{ex:vn-sigma-eig}

  \exercise{1} Show that $HXH=Z$, $HZH=X$ and $HYH=-Y$, and interpret: $H$ exchanges the computational basis with the $\ket\pm$ basis. Is $S$ Hermitian? Is it involutory? Compute $S^2$ and $S^4$.

  \exercise{2} Compute the action of $Z$ and $S$ on $\ket+$. Which gate maps $\ket+$ to $\ket-$? Which maps $\ket+$ to $\ket{+\ii}$? What are the Bloch vectors involved?

  \exercise{2} Start in $\ket0$ and apply $H$, $Z$, $H$ in this order. What is the final state? Repeat with $H$, $S$, $S$, $H$. Explain both results with the identities $HZH=X$ and $S^2=Z$.

  \exercise{2} Show that $e^{-\ii\theta\,\vec n\cdot\vec\sigma/2}=\cos\frac\theta2\Id-\ii\sin\frac\theta2\,\vec n\cdot\vec\sigma$. Compute $R_x(\theta)\ket0$ for $R_x(\theta)=e^{-\ii\theta X/2}$ and the Bloch vector you obtain. Express $H,X,Z$ as such exponentials, up to a global phase (for instance $X=\ii\,e^{-\ii\pi X/2}$).\label{ex:rotation-formula}

  \exercise{2} Let $U(t_1,t_2)=\exp(-\ii H(t_2-t_1)/\hbar)$ with $H=H^\dagger$. Prove that $U$ is unitary, that $U(t_2,t_3)U(t_1,t_2)=U(t_1,t_3)$, and that $U(t_2,t_1)=U(t_1,t_2)^{-1}$.

  \exercise{2} Prove that $K=-\ii\log U$ is Hermitian for every unitary $U$ and $U=e^{\ii K}$. Is $K$ unique? Find two different Hermitian $K$ with $e^{\ii K}=\Id$.

  \exercise{2} A qubit has $H=\hbar\omega Y$ and starts in $\ket0$. Show that $\ket{\psi(t)}=\cos(\omega t)\ket0+\sin(\omega t)\ket1$, find its Bloch vector, and describe the motion. Which state is the ground state?

  \exercise{3} Ehrenfest: for time-independent $A$, prove $\frac{d}{dt}\langle A\rangle_t=\frac{\ii}{\hbar}\langle[H,A]\rangle_t$. Deduce that $\langle H\rangle_t$ is constant and that every $A$ commuting with $H$ is a conserved quantity.

  \exercise{2} Let $H=\sum_EE\ket E\!\bra E$ and $\ket{\psi(0)}=\sum_Ec_E\ket E$. Show $\ket{\psi(t)}=\sum_Ec_Ee^{-\ii Et/\hbar}\ket E$ and that the probabilities $\abs{c_E}^2$ of the energy do not depend on time. For another observable $A$, show $\langle A\rangle_t=\sum_{E,E'}\bar c_Ec_{E'}A_{EE'}e^{\ii(E-E')t/\hbar}$, which oscillates at the Bohr frequencies $(E-E')/\hbar$.

  \exercise{2} Prove that commuting Hermitian operators have a common orthonormal eigenbasis (Proposition~\ref{prop:commuting-exp}). (Hint: the eigenspaces of $A$ are invariant under $B$.)\label{ex:commuting-basis}

  \exercise{2} Let $H(t)=\Omega(t)X$. Show $[H(t),H(t')]=0$ and $U=\exp\bigl(-\ii X\int_0^t\Omega(s)\dd s\bigr)$. Starting from $\ket0$, compute the probability of $\ket1$. If $\int_0^T\Omega=\pi/2$, which gate have you implemented?

  \exercise{2} Show that $e^{\ii s\sx}e^{\ii s\sz}\neq e^{\ii s(\sx+\sz)}$ for small $s\neq0$ by comparing the terms of order $s^2$. Then repeat Example~\ref{ex:trotter} numerically for $n=1,2,10$ and describe the error as $n$ grows.\label{ex:trotter-exercise}

  \exercise{3} How many real parameters describe a normalised state of $n$ qubits up to a global phase? Compare with the $2n$ parameters of a product of $n$ single-qubit states, and discuss what the gap suggests (see Section~\ref{sec:composite}).

  \exercise{1} Show that $U(g)=X$ and $U(g)=Z$ and $U(g)=H$ each define a unitary representation of $\Z_2=\{e,g\}$ on $\C^2$. Is $U(g)=S$ one? What group would $S$ represent?

  \exercise{2} Show that $U(g^{-1})=U(g)^\dagger$ in any unitary representation, and that if $G$ is abelian then all $U(g)$ commute.

  \exercise{2} Let $S\ket j=\ket{j+1\bmod N}$ on $\C^N$ and $\ket{f_k}=\frac1{\sqrt N}\sum_j\omega^{jk}\ket j$ with $\omega=e^{2\pi\ii/N}$. Show $S\ket{f_k}=\omega^{-k}\ket{f_k}$, so the Fourier basis diagonalises the representation $k\mapsto S^k$ of $\Z_N$.\label{ex:shift-fourier}

  \exercise{2} Verify that $\mathrm{SWAP}$ is Hermitian, unitary, with eigenvalues $\pm1$, and that $\tfrac12(\Id\pm\mathrm{SWAP})$ are the projectors onto the symmetric and antisymmetric subspaces. Which of the four Bell states lie in each?

  \exercise{2} Show that $H=\hbar\omega(Z\otimes\Id+\Id\otimes Z)$ commutes with $e^{\ii\varphi(Z\otimes\Id+\Id\otimes Z)}$ and with $\mathrm{SWAP}$, and use Proposition~\ref{prop:sym-energy} to explain the degeneracy of its energy levels.

  \exercise{3} Show that $R_{\vec n}(\theta)R_{\vec n}(\theta')=R_{\vec n}(\theta+\theta')$ and $R_{\vec n}(2\pi)=-\Id$. Deduce that $\theta\mapsto R_{\vec n}(\theta)$ is a representation of $\R$, and that it is a representation of the circle group $\R/4\pi\Z$ but not of $\R/2\pi\Z$.
\end{exerciseset}

%======================================================================
\section{Quantum measurement}\label{sec:measurement}
%======================================================================

\begin{objectives}
  You will state Postulate~3 and compute outcome probabilities and post-measurement states; prove that non-orthogonal states cannot be perfectly distinguished; work with projective measurements, expectation values and the uncertainty principle; explain why global phases are unobservable and relative phases are not;
  recognise positive operators; and use POVMs when only statistics matter.
\end{objectives}

\begin{recall}
  Orthogonal projectors and the spectral decomposition $M=\sum_mm\,P_m$ (Section~\ref{sec:operators}); the Cauchy--Schwarz inequality (Section~\ref{sec:hilbert}); Bloch vectors and the gates (Section~\ref{sec:state-evolution}).
\end{recall}

Unitary evolution describes a system that is left alone. But to learn anything about a system an experimentalist, who is an external system, must interact with it, and the system is then no longer closed. Postulate~3 describes what this does.

\subsection{General measurements}

\begin{postulate}[Measurement]\label{post:measurement}
  Quantum measurements are described by a collection $\{M_m\}$ of \term{measurement operators} acting on the state space. The index $m$ labels the possible outcomes.
  If the state immediately before the measurement is $\ket\psi$, then outcome $m$ occurs with probability and leaves the system in the state
  \[
    p(m)=\bra\psi M_m^\dagger M_m\ket\psi,\qquad
    \ket{\psi'_m}=\frac{M_m\ket\psi}{\sqrt{\bra\psi M_m^\dagger M_m\ket\psi}} .
  \]
  The operators satisfy the \term{completeness equation} $\sum_mM_m^\dagger M_m=\Id$.
\end{postulate}

Read the postulate as a recipe with three steps:
\begin{steps}
  \item apply each $M_m$ to the state; the squared length of $M_m\ket\psi$ is the probability of outcome $m$;
  \item the outcome that actually occurs is random, with these probabilities;
  \item the state after the measurement is $M_m\ket\psi$ rescaled to length $1$.
\end{steps}
Completeness is exactly what makes probabilities add to one:
\[
  \sum_mp(m)=\bra\psi\Bigl(\sum_mM_m^\dagger M_m\Bigr)\ket\psi=\braket\psi\psi=1\quad\text{for every state.}
\]
(Over $\C$ this requirement is equivalent to the operator identity, because if $\bra\psi A\ket\psi=0$ for all $\ket\psi$ then $A=0$, and the identity is much easier to check.)

\begin{example}\label{ex:meas-comp}
  Measure a qubit $\ket\psi=a\ket0+b\ket1$ in the computational basis, with $M_0=\ket0\!\bra0$ and $M_1=\ket1\!\bra1$.
\end{example}
\begin{solution}
  \begin{steps}
    \item \emph{Completeness.} Both operators are Hermitian with $M_m^2=M_m$, so $M_0^\dagger M_0+M_1^\dagger M_1=M_0+M_1=\Id$.
    \item \emph{Probabilities.} $p(0)=\bra\psi M_0\ket\psi=\abs a^2$ and $p(1)=\abs b^2$.
    \item \emph{Post-measurement states.} $\dfrac{M_0\ket\psi}{\abs a}=\dfrac a{\abs a}\ket0$ and $\dfrac{M_1\ket\psi}{\abs b}=\dfrac b{\abs b}\ket1$.
    \item \emph{Reading.} The factors $a/\abs a$ and $b/\abs b$ have modulus one: they are global phases, irrelevant (Section~\ref{sec:phase-global}). So the state \emph{collapses} to $\ket0$ or $\ket1$ with probabilities $\abs a^2,\abs b^2$. The amplitudes of Postulate~1 are, in this sense, the square roots of probabilities, which is why $\abs a^2+\abs b^2=1$ (Figure~\ref{fig:meas-qubit}).
  \end{steps}
\end{solution}

\begin{figure}[htb]
  \centering
  \begin{tikzpicture}[>=stealth,scale=1.3]
    \draw[->,ink!70] (-0.4,0) -- (3.0,0) node[right,font=\small] {$\ket0$};
    \draw[->,ink!70] (0,-0.4) -- (0,3.0) node[above,font=\small] {$\ket1$};
    \draw[ink!30] (0,0) circle (2.4);
    \draw[very thick,navy,->] (0,0) -- (35:2.4) node[above right] {$\ket\psi$};
    \draw[dashed,ink!60] (35:2.4) -- ({2.4*cos(35)},0) node[below,font=\small] {$a$};
    \draw[dashed,ink!60] (35:2.4) -- (0,{2.4*sin(35)}) node[left,font=\small] {$b$};
    \draw[very thick,figacc,->] (0,0) -- ({2.4*cos(35)},0);
    \draw[very thick,figacc,->] (0,0) -- (0,{2.4*sin(35)});
  \end{tikzpicture}
  \caption{For a state with real amplitudes, $\ket\psi=a\ket0+b\ket1$ with $a=\cos\alpha$, $b=\sin\alpha$. The outcome probabilities are the squared lengths $a^2,b^2$ of the two projections, and they add to $1$ by Pythagoras.}
  \label{fig:meas-qubit}
\end{figure}

\begin{example}[The same state in two bases]\label{ex:two-bases}
  Let $\ket\psi=\tfrac15\bigl(3\ket0+4\ii\ket1\bigr)$ (Example~\ref{ex:read-amplitudes}). Compare the outcome probabilities of the measurement in the computational basis and in the $\ket\pm$ basis.
\end{example}
\begin{solution}
  \begin{steps}
    \item \emph{Computational basis.} $p(0)=\tfrac9{25}=0.36$, $p(1)=\tfrac{16}{25}=0.64$.
    \item \emph{$\ket\pm$ basis: the overlaps.} $\braket\pm\psi=\tfrac1{\sqrt2}\cdot\tfrac15(3\pm4\ii)$.
    \item \emph{Squared moduli.} $\abs{3\pm4\ii}^2=25$, so $p(\pm)=\tfrac1{2}\cdot\tfrac{25}{25}=\tfrac12$ each.
  \end{steps}
  The state looks lopsided in one basis and perfectly balanced in the other. Probabilities depend on \emph{which question} we ask, and a state is a complete list of the answers to all of them.
\end{solution}

Measurement differs from evolution in two ways. Evolution is deterministic and reversible, $U^{-1}=U^\dagger$. Measurement is probabilistic, and typically irreversible: $M_m$ generally has a non-trivial kernel, so it loses information.

\begin{remark}
  A measuring device is itself a quantum system, so the system plus the apparatus form a larger isolated system, which by Postulate~2 evolves unitarily. Could Postulate~3 be \emph{derived} from Postulate~2? Despite a great deal of work, physicists still disagree; this is the \emph{measurement problem}. We take the pragmatic view: in practice it is clear when to apply Postulate~2 and when to apply Postulate~3.
\end{remark}

Measurements can be chained. Measuring with $\{L_l\}$ and then with $\{M_m\}$ turns out to be one single measurement.

\begin{proposition}[Cascades are single measurements]\label{prop:cascade}
  Measuring with $\{L_l\}$ and afterwards with $\{M_m\}$ is equivalent to a single measurement with operators $N_{lm}=M_mL_l$.
\end{proposition}
\begin{proof}
  \begin{steps}
    \item \emph{First measurement.} Outcome $l$ has probability $\bra\psi L_l^\dagger L_l\ket\psi$ and leaves $\ket{\psi_l}=L_l\ket\psi/\norm{L_l\psi}$.
    \item \emph{Second measurement.} Given this, outcome $m$ has conditional probability $\bra{\psi_l}M_m^\dagger M_m\ket{\psi_l}$.
    \item \emph{Joint probability.} Multiplying, $\PP(l,m)=\bra\psi L_l^\dagger M_m^\dagger M_mL_l\ket\psi=\bra\psi N_{lm}^\dagger N_{lm}\ket\psi$, and the final state is $N_{lm}\ket\psi$ normalised.
    \item \emph{Completeness.} $\sum_{l,m}N_{lm}^\dagger N_{lm}=\sum_lL_l^\dagger\bigl(\sum_mM_m^\dagger M_m\bigr)L_l=\sum_lL_l^\dagger L_l=\Id$. \qedhere
  \end{steps}
\end{proof}

Note the order: $N_{lm}=M_mL_l$, the first measurement acts first.

\begin{example}\label{ex:cascade}
  Measure first in the computational basis, $L_l=\ket l\!\bra l$, and then in the $\ket\pm$ basis, $M_\pm=\ket\pm\!\bra\pm$. Find the probabilities for $\ket\psi=a\ket0+b\ket1$.
\end{example}
\begin{solution}
  \begin{steps}
    \item \emph{The operators.} $N_{l\pm}=\ket\pm\!\bra\pm\ket l\!\bra l=\braket\pm l\,\ket\pm\!\bra l=\frac{(\pm1)^l}{\sqrt2}\ket\pm\!\bra l$.
    \item \emph{Acting on $\ket\psi$.} $N_{l\pm}\ket\psi=\frac{(\pm1)^l}{\sqrt2}\psi_l\ket\pm$, with $\psi_0=a$, $\psi_1=b$.
    \item \emph{Probabilities.} $\PP(l,\pm)=\tfrac12\abs{\psi_l}^2$, for each sign. The final state is $\ket\pm$ up to a phase.
    \item \emph{Completeness check.} $\sum_{l,\pm}\tfrac12\abs{\psi_l}^2=\abs a^2+\abs b^2=1$.
  \end{steps}
  The second measurement is completely random whatever the first outcome: after having been forced into $\ket0$ or $\ket1$, the state has no memory of its previous superposition.
\end{solution}

\subsection{Distinguishing quantum states}

Here is a game. Alice chooses a state $\ket{\psi_i}$, $1\le i\le n$, from a set known to both players and gives it to Bob, who must name $i$. Classically, distinct states can be told apart, at least in principle. Quantum mechanically the answer depends on the geometry of the set.

\begin{proposition}\label{prop:distinguish-orthogonal}
  If the states $\ket{\psi_i}$ are orthonormal, Bob can identify $i$ with certainty.
\end{proposition}
\begin{proof}
  \begin{steps}
    \item \emph{The measurement.} Let Bob use $M_i=\ket{\psi_i}\!\bra{\psi_i}$ for $i\ge1$ and $M_0=\sqrt{\Id-\sum_{i\ge1}\ket{\psi_i}\!\bra{\psi_i}}$.
    \item \emph{Completeness.} The sum $\sum_{i\ge1}\ket{\psi_i}\!\bra{\psi_i}$ is a projector, so $\Id$ minus it is a projector and $M_0^\dagger M_0=\Id-\sum_{i\ge1}\ket{\psi_i}\!\bra{\psi_i}$. Hence $\sum_iM_i^\dagger M_i=\Id$.
    \item \emph{Success.} If $\ket{\psi_i}$ was prepared, $p(i)=\bra{\psi_i}M_i\ket{\psi_i}=1$. \qedhere
  \end{steps}
\end{proof}

\begin{theorem}[Non-orthogonal states cannot be distinguished]\label{thm:no-distinguish}
  If $\braket{\psi_1}{\psi_2}\neq0$, no measurement can identify with certainty which of $\ket{\psi_1},\ket{\psi_2}$ was prepared.
\end{theorem}
\begin{proof}
  Suppose, for contradiction, that Bob measures $\{M_j\}$ and guesses $i=f(j)\in\{1,2\}$ from the outcome $j$, and is always right.
  \begin{steps}
    \item \emph{Group outcomes by guess.} Define $E_i=\sum_{j:f(j)=i}M_j^\dagger M_j$, so $E_1+E_2=\Id$.
    \item \emph{``Always right''} means $\bra{\psi_1}E_1\ket{\psi_1}=1=\bra{\psi_2}E_2\ket{\psi_2}$.
    \item \emph{A consequence.} Since $E_1+E_2=\Id$, we get $\bra{\psi_1}E_2\ket{\psi_1}=0$, that is $\norm{\sqrt{E_2}\psi_1}=0$, so $\sqrt{E_2}\ket{\psi_1}=0$.
    \item \emph{Decompose $\ket{\psi_2}$} as in the proof of Cauchy--Schwarz: $\ket{\psi_2}=\alpha\ket{\psi_1}+\beta\ket\varphi$ with $\ket\varphi\perp\ket{\psi_1}$ a unit vector, $\alpha=\braket{\psi_1}{\psi_2}$, $\abs\alpha^2+\abs\beta^2=1$, and $\abs\beta<1$ because $\alpha\neq0$.
    \item \emph{Contradiction.} $\sqrt{E_2}\ket{\psi_2}=\beta\sqrt{E_2}\ket\varphi$, so $\bra{\psi_2}E_2\ket{\psi_2}=\abs\beta^2\bra\varphi E_2\ket\varphi\le\abs\beta^2<1$ (as $E_2\le\Id$). This contradicts $\bra{\psi_2}E_2\ket{\psi_2}=1$. \qedhere
  \end{steps}
\end{proof}

So orthogonality is the only thing that makes states perfectly distinguishable. In particular $\ket0$ and $\ket+$ cannot be told apart reliably. This is the basis of everything from quantum key distribution to the no-cloning theorem.

\subsection{Projective measurements}

An important special case is when the measurement operators are orthogonal projectors.

\begin{definition}[Projective measurement]\label{def:projective}
  A \term{projective measurement} is described by an \term{observable}, a Hermitian operator $M$ with spectral decomposition $M=\sum_mm\,P_m$, where $P_m$ is the projector onto the eigenspace of eigenvalue $m$. The outcomes are the eigenvalues, and
  \[
    p(m)=\bra\psi P_m\ket\psi,\qquad \ket{\psi'_m}=\frac{P_m\ket\psi}{\sqrt{p(m)}} .
  \]
\end{definition}

This is Postulate~3 with $M_m=P_m$: Hermitian with $P_mP_{m'}=\delta_{mm'}P_m$, so $\sum_mM_m^\dagger M_m=\sum_mP_m=\Id$. Two ways of specifying a projective measurement are equivalent: give the observable, or give the complete family of orthogonal projectors.
``Measure in the basis $\{\ket m\}$'' means $P_m=\ket m\!\bra m$. Projective measurements are \emph{repeatable}: measuring again right away gives the same $m$ and does not change the state, since $P_mP_m\ket\psi=P_m\ket\psi$.

The average outcome needs no eigenvalues:
\[
  \langle M\rangle=\sum_mm\,p(m)=\sum_mm\bra\psi P_m\ket\psi=\bra\psi M\ket\psi,
\]
and the \term{standard deviation} follows from $(\Delta M)^2=\langle(M-\langle M\rangle)^2\rangle=\langle M^2\rangle-\langle M\rangle^2$. Because $M-\langle M\rangle$ is Hermitian, $(\Delta M)^2=\norm{(M-\langle M\rangle)\psi}^2\ge0$, and $\Delta M=0$ exactly when $\ket\psi$ is an eigenvector of $M$.
If we prepare $\ket\psi$ many times and measure $M$ each time, $\langle M\rangle$ is the mean and $\Delta M$ the typical spread of the results (this is just the expectation and standard deviation of the random variable ``outcome'', see Section~\ref{sec:prob} and Example~\ref{ex:rv-qm}). A state has a definite value of $M$ iff it is an eigenstate of $M$.

\begin{example}\label{ex:meas-pauli}
  Let $\ket\psi=\frac1{\sqrt3}\bigl(\ket0+(1+\ii)\ket1\bigr)$. Find the outcomes, probabilities, $\langle M\rangle$ and $\Delta M$ for $M=X,Y,Z$.
\end{example}
\begin{solution}
  \begin{steps}
    \item \emph{Outcomes and projectors.} Each of $X,Y,Z$ has outcomes $\pm1$ with projectors $P_\pm=\tfrac12(\Id\pm\sigma_j)$.
    \item \emph{Probabilities.} By Example~\ref{ex:parseval}: for $Z$ they are $\tfrac13,\tfrac23$; for $X$ they are $\abs{\braket+\psi}^2=\tfrac56$ and $\abs{\braket-\psi}^2=\tfrac16$; for $Y$ they are $\abs{\braket{+\ii}\psi}^2=\tfrac56$ and $\tfrac16$.
    \item \emph{Expectations.} $\langle X\rangle=\tfrac56-\tfrac16=\tfrac23$, $\langle Y\rangle=\tfrac23$, $\langle Z\rangle=\tfrac13-\tfrac23=-\tfrac13$.
    \item \emph{Standard deviations.} Since $M^2=\Id$, $(\Delta M)^2=1-\langle M\rangle^2$, so $\Delta X=\Delta Y=\tfrac{\sqrt5}3$ and $\Delta Z=\tfrac{\sqrt8}3$.
    \item \emph{Check.} $(\tfrac23,\tfrac23,-\tfrac13)$ is the Bloch vector of $\ket\psi$ (Proposition~\ref{prop:bloch}), and $\abs{\vec n}=1$.
  \end{steps}
\end{solution}

For a real unit vector $\vec v$, the observable $\vec v\cdot\vec\sigma$ is called, for historical reasons, \emph{spin along the axis $\vec v$}. By Exercise~\ref{ex:vn-sigma-eig2} it has eigenvalues $\pm1$ with projectors $P_\pm=\tfrac12(\Id\pm\vec v\cdot\vec\sigma)$.
Measuring it on a state with Bloch vector $\vec n$ gives $+1$ with probability $\tfrac12(1+\vec n\cdot\vec v)$.

\begin{example}[Tilted measurement axis]\label{ex:tilted}
  A qubit is in $\ket0$. It is measured along $\vec v=(\sin\theta,0,\cos\theta)$. Find the probability of $+1$ and the post-measurement state when $+1$ occurs.
\end{example}
\begin{solution}
  \begin{steps}
    \item \emph{The Bloch vector of $\ket0$} is $\vec n=(0,0,1)$, so $\vec n\cdot\vec v=\cos\theta$.
    \item \emph{The probability.} $p(+1)=\tfrac12(1+\cos\theta)=\cos^2\tfrac\theta2$.
    \item \emph{The post-measurement state.} $P_+\ket0=\tfrac12(\Id+\vec v\cdot\vec\sigma)\ket0=\tfrac12\bigl((1+\cos\theta)\ket0+\sin\theta\ket1\bigr)=\cos\tfrac\theta2\bigl(\cos\tfrac\theta2\ket0+\sin\tfrac\theta2\ket1\bigr)$.
    \item \emph{Normalise.} Dividing by $\cos\frac\theta2=\sqrt{p(+1)}$ gives $\cos\tfrac\theta2\ket0+\sin\tfrac\theta2\ket1$, the state with Bloch vector $\vec v$.
  \end{steps}
  This is the quantum version of Malus's law for polarisers: the probability falls as $\cos^2$ of half the angle on the Bloch sphere, and a successful outcome leaves the system pointing exactly along the measurement axis.
\end{solution}

\subsection{The uncertainty principle}

If $C$ and $D$ do not commute, that is, if their commutator $[C,D]=CD-DC$ (Section~\ref{sec:vector-spaces}) is not zero, they cannot both be sharp in the same state. The precise statement is the following.

\begin{theorem}[Heisenberg--Robertson]\label{thm:uncertainty}
  For Hermitian $C,D$ and any state $\ket\psi$, $\ \Delta C\,\Delta D\ge\tfrac12\bigl|\langle[C,D]\rangle\bigr|$.
\end{theorem}
\begin{proof}
  \begin{steps}
    \item \emph{Centre the operators.} Let $A=C-\langle C\rangle$ and $B=D-\langle D\rangle$, both Hermitian, so $[A,B]=[C,D]$, $\langle A^2\rangle=(\Delta C)^2$ and $\langle B^2\rangle=(\Delta D)^2$.
    \item \emph{Real and imaginary parts.} Write $\langle AB\rangle=x+\ii y$ with $x,y\in\R$. Since $\langle BA\rangle=\overline{\langle AB\rangle}$, we get $\langle[A,B]\rangle=2\ii y$ and $\langle\{A,B\}\rangle=2x$. Hence $\abs{\langle[A,B]\rangle}^2+\abs{\langle\{A,B\}\rangle}^2=4\abs{\langle AB\rangle}^2$.
    \item \emph{Cauchy--Schwarz.} $\abs{\langle AB\rangle}^2=\abs{\braket{A\psi}{B\psi}}^2\le\langle A^2\rangle\langle B^2\rangle$.
    \item \emph{Drop the anticommutator term} (it is $\ge0$): $\abs{\langle[A,B]\rangle}^2\le4\langle A^2\rangle\langle B^2\rangle=4(\Delta C)^2(\Delta D)^2$.
    \item \emph{Square roots} give the theorem. \qedhere
  \end{steps}
\end{proof}

\begin{example}\label{ex:uncertainty-check}
  Check the uncertainty relation for $X,Y$ on the state of Example~\ref{ex:meas-pauli}, and on $\ket0$.
\end{example}
\begin{solution}
  \begin{steps}
    \item \emph{The right-hand side.} $[X,Y]=2\ii Z$, so $\tfrac12\abs{\langle[X,Y]\rangle}=\abs{\langle Z\rangle}$.
    \item \emph{The state of Example~\ref{ex:meas-pauli}.} $\Delta X\Delta Y=\tfrac59\ge\tfrac13=\abs{\langle Z\rangle}$.
    \item \emph{The state $\ket0$.} $\langle Z\rangle=1$ and $\Delta X=\Delta Y=1$ (outcomes $\pm1$ with probability $\tfrac12$ each), so $1\ge1$: equality.
  \end{steps}
\end{solution}

\begin{pitfall}
  The principle is often misquoted as ``measuring $C$ accurately disturbs $D$''. Measurement does disturb the system (Postulate~3), but that is not the content of the theorem. The statement is statistical: prepare many copies of $\ket\psi$, measure $C$ on some and $D$ on the others, and the spreads of the two lists obey the inequality.
  It is a property of the state and the operators, not of the quality of an apparatus. Commuting observables can have a common eigenbasis, hence states with $\Delta C=\Delta D=0$.
\end{pitfall}

For position and momentum, $[\hat x,\hat p]=\ii\hbar$ gives $\Delta x\,\Delta p\ge\hbar/2$; we saw in Example~\ref{ex:gaussian-uncertainty} that Gaussians attain this bound. Taking the trace of $[A,B]=\ii\hbar\Id$ in finite dimension would give $0=\ii\hbar n$, absurd (Exercise~\ref{ex:no-xp}): position and momentum need infinite-dimensional Hilbert spaces.

\subsection{Phase: global and relative}\label{sec:phase-global}

\begin{proposition}\label{prop:global-phase}
  For every $\gamma\in\R$, the states $\ket\psi$ and $e^{\ii\gamma}\ket\psi$ give identical statistics for every measurement, with post-measurement states that agree up to a global phase.
\end{proposition}
\begin{proof}
  \begin{steps}
    \item \emph{Probabilities.} $\bra\psi e^{-\ii\gamma}M_m^\dagger M_me^{\ii\gamma}\ket\psi=\bra\psi M_m^\dagger M_m\ket\psi$: the phase cancels.
    \item \emph{States.} $M_me^{\ii\gamma}\ket\psi=e^{\ii\gamma}M_m\ket\psi$: the phase is carried along. \qedhere
  \end{steps}
\end{proof}

Hence \emph{global phases are physically irrelevant}: a state is really a \emph{ray} $\{e^{\ii\gamma}\ket\psi\}$, the matrix $\ket\psi\!\bra\psi$ is blind to $\gamma$, and for a qubit the set of rays is the Bloch sphere. \emph{Relative} phases, on the other hand, matter. The states $\ket\pm=\frac1{\sqrt2}(\ket0\pm\ket1)$ have the same amplitude moduli, $\frac1{\sqrt2}$ for $\ket0$, and amplitudes for $\ket1$ differing by the factor $-1$.
Measured in the computational basis, they give the same statistics $(\tfrac12,\tfrac12)$; but measured in the $\ket\pm$ basis, $\ket+$ gives $+1$ with certainty and $\ket-$ gives $-1$ with certainty. States differing by a relative phase are physically different, and the relative phase is basis dependent. This is exactly the phenomenon of Example~\ref{ex:relative-phase}.

\begin{example}\label{ex:phase-pairs}
  Decide which of these pairs describe the same physical state: (a) $\ket\psi$ and $-\ket\psi$; (b) $\frac1{\sqrt2}(\ket0+\ket1)$ and $\frac1{\sqrt2}(\ket0+\ii\ket1)$; (c) $\frac1{\sqrt2}(\ket0+\ii\ket1)$ and $\frac1{\sqrt2}(\ii\ket0-\ket1)$. For the pairs that differ, find a measurement that distinguishes them.
\end{example}
\begin{solution}
  \begin{steps}
    \item (a) Same: $-1=e^{\ii\pi}$ is a global phase.
    \item (b) Different: they are $\ket+$ and $\ket{+\ii}$, with Bloch vectors $(1,0,0)$ and $(0,1,0)$, and no global phase relates them. Measuring $X$ gives $+1$ with probability $1$ on the first and $\tfrac12$ on the second.
    \item (c) Same: $\frac1{\sqrt2}(\ii\ket0-\ket1)=\ii\cdot\frac1{\sqrt2}(\ket0+\ii\ket1)$, because $\ii\cdot\ii=-1$.
  \end{steps}
\end{solution}

\begin{example}\label{ex:phase-xy}
  For $\ket{\psi_\varphi}=\frac1{\sqrt2}(\ket0+e^{\ii\varphi}\ket1)$, find the probabilities of the outcome $+1$ when measuring $X$ and $Y$.
\end{example}
\begin{solution}
  \begin{steps}
    \item \emph{The Bloch vector} is $\vec n=(\cos\varphi,\sin\varphi,0)$ (here $\theta=\pi/2$).
    \item \emph{The probabilities.} By the rule above, $P(X=+1)=\tfrac12(1+\cos\varphi)$ and $P(Y=+1)=\tfrac12(1+\sin\varphi)$.
    \item \emph{Reading.} As $\varphi$ varies, both probabilities oscillate; the relative phase is directly observable.
  \end{steps}
\end{solution}

\subsection{Positive operators}

Before going on we need a concept that will be used for measurements and for states alike: that of a positive operator. It is the operator analogue of a non-negative number.

\begin{definition}[Positive operator]\label{def:positive}
  An operator $A$ on $\Hil$ is \term{positive} (written $A\ge0$) if $\bra\psi A\ket\psi\ge0$ for every $\ket\psi\in\Hil$. We write $A\ge B$ if $A-B\ge0$ (the \term{L\"owner order}).
\end{definition}

The definition requires $\bra\psi A\ket\psi$ to be real and non-negative for all $\ket\psi$; over $\C$ the realness forces $A$ to be Hermitian. The following criteria are how we recognise positive operators in practice.

\begin{theorem}[Characterisations of positivity]\label{thm:positive}
  For an operator $A$ on a finite-dimensional Hilbert space, the following are equivalent.
  \begin{enumerate}
    \item[(i)] $A\ge0$.
    \item[(ii)] $A$ is Hermitian with all eigenvalues $\ge0$.
    \item[(iii)] $A=C^2$ for some positive operator $C$ (one writes $C=\sqrt A$, it is unique).
    \item[(iv)] $A=B^\dagger B$ for some operator $B$.
  \end{enumerate}
\end{theorem}
\begin{proof}
  \begin{steps}
    \item \emph{(i)$\Rightarrow$(ii).} Write $A=B+\ii C$ with $B=\tfrac12(A+A^\dagger)$ and $C=\tfrac1{2\ii}(A-A^\dagger)$ both Hermitian. Then $\bra\psi A\ket\psi=\bra\psi B\ket\psi+\ii\bra\psi C\ket\psi$ with both expectation values real, so reality of the sum forces $\bra\psi C\ket\psi=0$ for all $\ket\psi$, hence $C=0$ and $A=B$ is Hermitian. If $A\ket v=\lambda\ket v$ with $\norm v=1$, then $\lambda=\bra vA\ket v\ge0$.
    \item \emph{(ii)$\Rightarrow$(iii).} By the spectral theorem $A=\sum_i\lambda_i\ket{v_i}\!\bra{v_i}$ with $\lambda_i\ge0$. Put $C=\sum_i\sqrt{\lambda_i}\ket{v_i}\!\bra{v_i}$, which is Hermitian with eigenvalues $\ge0$ (so $C\ge0$ by the direction (ii)$\Rightarrow$(i), below), and $C^2=A$.
    \item \emph{(iii)$\Rightarrow$(iv).} Take $B=C$; then $B^\dagger B=C^2=A$ since $C$ is Hermitian.
    \item \emph{(iv)$\Rightarrow$(i).} $\bra\psi B^\dagger B\ket\psi=\norm{B\psi}^2\ge0$.
    \item \emph{(ii)$\Rightarrow$(i).} If $A=\sum_i\lambda_i\ket{v_i}\!\bra{v_i}$ with $\lambda_i\ge0$ then $\bra\psi A\ket\psi=\sum_i\lambda_i\abs{\braket{v_i}\psi}^2\ge0$. \qedhere
  \end{steps}
\end{proof}

For $2\times2$ matrices there is a quick test: a Hermitian matrix $\begin{pmatrix}a&b\\\bar b&d\end{pmatrix}$ is positive iff $a\ge0$, $d\ge0$ and $ad-\abs b^2\ge0$ (trace and determinant are the sum and the product of the eigenvalues, and two real numbers with non-negative sum and product are both non-negative).

\begin{example}[Recognising positive operators]\label{ex:positive-examples}
  Decide which are positive: (a) $\ket\psi\!\bra\psi$ and any projector; (b) $\begin{pmatrix}2&1\\1&2\end{pmatrix}$; (c) $\begin{pmatrix}1&2\\2&1\end{pmatrix}$; (d) $Z$ and $\Id+Z$; (e) $A^\dagger A$ for $A=\begin{pmatrix}1&1\\0&1\end{pmatrix}$.
\end{example}
\begin{solution}
  \begin{steps}
    \item (a) $\bra\varphi\bigl(\ket\psi\!\bra\psi\bigr)\ket\varphi=\braket\varphi\psi\braket\psi\varphi=\abs{\braket\psi\varphi}^2\ge0$: positive. A projector has eigenvalues $0,1$: positive.
    \item (b) Hermitian with eigenvalues $3$ and $1$ (Example~\ref{ex:eigen}): positive. (Test: $a=d=2$, $ad-\abs b^2=3\ge0$.)
    \item (c) Hermitian, but $ad-\abs b^2=1-4<0$: the eigenvalues are $3$ and $-1$. \emph{Not} positive, although all entries are positive numbers.
    \item (d) $Z$ has eigenvalue $-1$: not positive. $\Id+Z=2\ket0\!\bra0$ has eigenvalues $2,0$: positive.
    \item (e) $A^\dagger A=\begin{pmatrix}1&1\\1&2\end{pmatrix}$, which is of the form $B^\dagger B$, so positive by (iv). Check: $ad-\abs b^2=2-1=1\ge0$.
  \end{steps}
\end{solution}

\begin{pitfall}
  ``Positive'' does not mean ``has positive entries'': the matrix in (c) has only positive entries and is not a positive operator. The test is on the \emph{eigenvalues} (or on $\bra\psi A\ket\psi$ for all $\ket\psi$), which does not depend on the basis.
\end{pitfall}

\begin{lookahead}
  Positivity is the right notion of ``cannot give a negative probability'' for operators, and it will appear twice in this chapter. In the next subsection the operators $E_m=M_m^\dagger M_m$ giving measurement probabilities are positive by (iv), and form a POVM. In Section~\ref{sec:density} the density operators will be characterised as the positive operators of trace one, and the identity $\rho=(\Id+\vec r\cdot\vec\sigma)/2\ge0\iff\abs{\vec r}\le1$ is exactly the positivity test above applied to a qubit.
\end{lookahead}

\subsection{POVM measurements}

Postulate~3 gives two things: the statistics $p(m)$ and the post-measurement state. Often only the statistics matter, for instance when the system is measured once and then discarded. In that case it is enough to know the operators $E_m=M_m^\dagger M_m$, since $p(m)=\bra\psi E_m\ket\psi$.

\begin{definition}[POVM]\label{def:povm}
  A \term{POVM} (positive operator-valued measure) is a family $\{E_m\}$ of operators such that (a) each $E_m$ is positive; (b) $\sum_mE_m=\Id$.
\end{definition}

Every measurement $\{M_m\}$ defines a POVM $E_m=M_m^\dagger M_m$, positive by Theorem~\ref{thm:positive}(iv). Conversely, every POVM arises this way: take $M_m=\sqrt{E_m}$ (Theorem~\ref{thm:positive}(iii)), so $\sum_mM_m^\dagger M_m=\sum_mE_m=\Id$. The POVM determines each $M_m$ only up to a unitary applied afterwards: by the polar decomposition,
every operator can be written $M=U\sqrt{M^\dagger M}$ with $U$ unitary (the operator analogue of $z=e^{\ii\theta}\abs z$). Two measurements with the same POVM have identical statistics and differ only by what the apparatus does afterwards. In the projective case, and only then, $E_m=P_m^\dagger P_m=P_m$.

\begin{example}[Never wrong, sometimes silent]\label{ex:unambiguous}
  Alice gives Bob a qubit in $\ket{\psi_1}=\ket0$ or $\ket{\psi_2}=\ket+$. Since these are not orthogonal, Bob cannot always decide (Theorem~\ref{thm:no-distinguish}). Show that he can nevertheless decide \emph{some} of the time without ever being wrong.
\end{example}
\begin{solution}
  \begin{steps}
    \item \emph{The idea.} Make outcome $1$ impossible for $\ket0$ (so it can only come from $\ket+$) and outcome $2$ impossible for $\ket+$. The operator that annihilates $\ket0$ is $\ket1\!\bra1$; the one that annihilates $\ket+$ is $\ket-\!\bra-$.
    \item \emph{The POVM.} With $c=\frac{\sqrt2}{1+\sqrt2}=2-\sqrt2\approx0.586$, put $E_1=c\ket1\!\bra1$, $E_2=c\ket-\!\bra-$, $E_3=\Id-E_1-E_2$.
    \item \emph{Positivity.} $E_1,E_2$ are positive. For $E_3$: $\ket1\!\bra1+\ket-\!\bra-$ has eigenvalues $1\pm\abs{\braket1-}=1\pm\tfrac1{\sqrt2}$, so $E_3$ has eigenvalues $1-c(1\pm\tfrac1{\sqrt2})$, which are $0$ and $\approx0.83$: both $\ge0$.
    \item \emph{Never wrong.} $E_1\ket0=0$, so outcome $1$ is impossible for $\ket0$: if it occurs, the state was $\ket+$. Likewise $E_2\ket+=0$.
    \item \emph{Probabilities.}
      \begin{center}
      \begin{tabular}{lccc}
        \toprule
        input & $p(E_1)$ & $p(E_2)$ & $p(E_3)$\\
        \midrule
        $\ket0$ & $0$ & $c\abs{\braket-0}^2=c/2$ & $1-c/2$\\
        $\ket+$ & $c\abs{\braket1+}^2=c/2$ & $0$ & $1-c/2$\\
        \bottomrule
      \end{tabular}
      \end{center}
    \item \emph{The price.} With probability $1-c/2\approx0.71$ Bob gets $E_3$ and learns nothing. The success probability is $c/2=1-\tfrac1{\sqrt2}=1-\abs{\braket0+}\approx0.29$.
  \end{steps}
\end{solution}

We can summarise the three formalisms.
\begin{center}
\small
\begin{tabular}{@{}p{0.13\textwidth}p{0.17\textwidth}p{0.33\textwidth}p{0.13\textwidth}p{0.12\textwidth}@{}}
  \toprule
  Type & Data & Conditions & Post-state & Repeatable\\
  \midrule
  General & $\{M_m\}$ & $\sum M_m^\dagger M_m=\Id$ & yes & not in general\\
  Projective & $M=\sum mP_m$ & $P_m^\dagger=P_m$, $P_mP_{m'}=\delta_{mm'}P_m$, $\sum P_m=\Id$ & yes & yes\\
  POVM & $\{E_m\}$ & $E_m\ge0$, $\sum E_m=\Id$ & no (statistics only) & ---\\
  \bottomrule
\end{tabular}
\end{center}
Why bother with the general formalism? It is simpler mathematically, some problems have optimal solutions that are not projective (optimal state discrimination, the trine POVM of Exercise~\ref{ex:trine}), and many real measurements are not repeatable: a screen that records a photon destroys it.
The good news, shown in Section~\ref{sec:composite}, is that once we have unitaries and composite systems, projective measurements suffice: ancillas turn them into general measurements.

\begin{exerciseset}
  \exercise{1} Suppose the system is in an eigenstate $\ket\psi$ of an observable $M$ with eigenvalue $m$. What are the average observed value and the standard deviation?

  \exercise{1} A qubit is in $\ket0$ and we measure $X$. What are $\langle X\rangle$ and $\Delta X$? Which outcomes can occur, with which probabilities, and what are the post-measurement states?

  \exercise{2} Show $(\vec v\cdot\vec\sigma)^2=\Id$ for a real unit vector $\vec v$ (use $\sigma_j\sigma_k=\delta_{jk}\Id+\ii\sum_l\varepsilon_{jkl}\sigma_l$), deduce that its eigenvalues are $\pm1$ and that the projectors are $P_\pm=\frac12(\Id\pm\vec v\cdot\vec\sigma)$. Compute the probability of the result $+1$ when the state is $\ket0$, and the state afterwards.\label{ex:vn-sigma-eig2}

  \exercise{2} Let $\ket\psi=\frac1{\sqrt3}\ket0+\frac{1+\ii}{\sqrt3}\ket1$. Verify Example~\ref{ex:meas-pauli}, and check the three uncertainty relations $\Delta X\Delta Y\ge\abs{\langle Z\rangle}$ and its two cyclic versions.

  \exercise{2} For the state $\ket\psi=\tfrac15(3\ket0+4\ii\ket1)$ compute $\langle X\rangle,\langle Y\rangle,\langle Z\rangle$, the three standard deviations, and the Bloch vector; check the uncertainty relation for each pair.

  \exercise{2} Cascades: check the four operators $N_{l\pm}=M_\pm L_l$ of Example~\ref{ex:cascade} satisfy $\sum N^\dagger N=\Id$, and compute probabilities and final states for $\ket\psi=a\ket0+b\ket1$. Show the order matters by repeating with $\{M_\pm\}$ first and $\{L_l\}$ second.

  \exercise{2}\label{ex:epsilon-meas} For $0\le\epsilon\le1$ let $M_0=\ket0\!\bra0+\sqrt{1-\epsilon}\ket1\!\bra1$ and $M_1=\sqrt\epsilon\ket1\!\bra1$. Check completeness, compute $p(0),p(1)$ and the post-measurement states for $\ket\psi=a\ket0+b\ket1$. What happens for $\epsilon=0$ and $\epsilon=1$? Show that for $0<\epsilon<1$ the measurement is neither projective nor repeatable, and find its POVM.

  \exercise{1} Show that $\ket0,\ket1$ can be distinguished perfectly but $\ket0,\ket+$ cannot. What is the error probability of the best projective guess for the latter (equal priors)? Compare with the unambiguous POVM of Example~\ref{ex:unambiguous}.

  \exercise{2} Let $f$ be a real function on the spectrum of $M$. Show that $\mathbb E[f(M)]=\bra\psi f(M)\ket\psi$, and verify it for $f(M)=M^2$ with $M=Z$.

  \exercise{3} Prove that $\Delta M=0$ iff $\ket\psi$ is an eigenvector of $M$, and that $(\Delta M)^2=\langle M^2\rangle-\langle M\rangle^2$.

  \exercise{3} For $\ket\psi=\cos\frac\theta2\ket0+e^{\ii\phi}\sin\frac\theta2\ket1$ with Bloch vector $\vec n$, show $\langle X\rangle=n_x$, $\langle Y\rangle=n_y$, $\langle Z\rangle=n_z$ and $(\Delta X)^2=1-n_x^2$, $(\Delta Y)^2=1-n_y^2$. Show that $(\Delta X)^2(\Delta Y)^2=n_z^2+n_x^2n_y^2$, check the uncertainty relation, and find when equality holds.

  \exercise{2}\label{ex:no-xp} On smooth functions vanishing at infinity let $\hat xf=xf$ and $\hat pf=-\ii\hbar f'$. Show that $[\hat x,\hat p]=\ii\hbar\Id$. Using the trace, show that no finite matrices $A,B$ satisfy $[A,B]=\ii\Id$.

  \exercise{2} Express $\ket\pm$ in a basis in which they are \emph{not} the same up to a relative phase. Decide which of the following pairs are the same physical state, and for those that are not find a projective measurement whose statistics differ: $\ket\psi$ and $-\ket\psi$; $\frac1{\sqrt2}(\ket0+\ket1)$ and $\frac1{\sqrt2}(\ket0+\ii\ket1)$; $\frac1{\sqrt2}(\ket0+\ii\ket1)$ and $\frac1{\sqrt2}(\ii\ket0-\ket1)$.

  \exercise{2} For $\ket{\psi_\varphi}=\frac1{\sqrt2}(\ket0+e^{\ii\varphi}\ket1)$ compute the probabilities of the outcome $+1$ for $X$ and $Y$ as functions of $\varphi$. Prove that two states whose amplitudes have the same moduli in the computational basis cannot be distinguished by a measurement in that basis, and give an explicit measurement that distinguishes them when they differ by a non-trivial relative phase.

  \exercise{1} Decide which are positive: $\begin{pmatrix}1&0\\0&0\end{pmatrix}$, $\begin{pmatrix}0&1\\1&0\end{pmatrix}$, $\begin{pmatrix}3&\ii\\-\ii&1\end{pmatrix}$, $\begin{pmatrix}1&2\\2&3\end{pmatrix}$, $X+2\Id$, $\ket+\!\bra+-\ket-\!\bra-$.

  \exercise{2} Prove the $2\times2$ test: $\begin{pmatrix}a&b\\\bar b&d\end{pmatrix}\ge0$ iff $a\ge0$, $d\ge0$ and $ad\ge\abs b^2$. Find all real $c$ such that $\begin{pmatrix}1&c\\c&4\end{pmatrix}\ge0$.

  \exercise{2} Find $\sqrt A$ for $A=\begin{pmatrix}5&-3\\-3&5\end{pmatrix}$ and for $A=\ket\psi\!\bra\psi$ with $\norm\psi=1$.

  \exercise{2} Show that if $A,B\ge0$ then $A+B\ge0$ and $\Tr A\ge0$, and that $CAC^\dagger\ge0$ for every $C$. Is $AB\ge0$ in general? (Try $A=\ket0\!\bra0$, $B=\ket+\!\bra+$.)

  \exercise{3} Show that if $A,B\ge0$ then $\Tr(AB)\ge0$. (Hint: write $B=\sum_j\mu_j\ket{w_j}\!\bra{w_j}$.)

  \exercise{2} Prove that $\sqrt{A}$ is unique for $A\ge0$ when $A$ has distinct eigenvalues. Show that $-\sqrt A$ is a square root of $A$ that is not positive.

  \exercise{2} Verify that $E_1,E_2$ of Example~\ref{ex:unambiguous} are positive and that $E_3=\Id-E_1-E_2$ is positive (compute the eigenvalues of $\ket1\!\bra1+\ket-\!\bra-$). For two states with overlap $s=\abs{\braket{\psi_1}{\psi_2}}$, conjecture and test the optimal success probability of unambiguous discrimination (it equals $1-s$; check $s=0$ and $s=1$).

  \exercise{2}\label{ex:trine} (Trine.) Let $\vec n_k=(\sin\theta_k,0,\cos\theta_k)$ with $\theta_k=2\pi k/3$ and $P_k=\frac12(\Id+\vec n_k\cdot\vec\sigma)$ for $k=0,1,2$. Show $E_k=\frac23P_k$ is a POVM (hint: $\sum_k\vec n_k=0$), and compute $p(k)$ for the state $\ket0$. Explain why it is not projective.

  \exercise{2} Show that any measurement in which $M_m=E_m=M_m^\dagger M_m$ is a projective measurement. (Hint: first $M_m^\dagger=M_m$ and $M_m^2=M_m$; then use $\sum_mM_m=\Id$ to show $M_mM_{m'}=0$ for $m\neq m'$.)

  \exercise{3} Given linearly independent states $\ket{\psi_1},\dots,\ket{\psi_m}$, construct a POVM $\{E_1,\dots,E_{m+1}\}$ such that outcome $E_i$ ($i\le m$) tells Bob with certainty that the state was $\ket{\psi_i}$, with $\bra{\psi_i}E_i\ket{\psi_i}>0$. (Hint: use dual-basis vectors $\ket{\psi_i^*}\perp\ket{\psi_j}$ for $j\neq i$, and $E_i=c\ket{\psi_i^*}\!\bra{\psi_i^*}$ with $c>0$ small.)
\end{exerciseset}

%======================================================================
\section{Composite systems and entanglement}\label{sec:composite}
%======================================================================

\begin{objectives}
  You will build the state space of several systems with the tensor product and compute with Kronecker products;
  decide whether a two-qubit state is entangled; measure one part of an entangled system;
  and see how general measurements arise from projective ones on a larger system.
\end{objectives}

A single qubit is not very interesting. To compute with qubits we need many of them, and to describe several quantum systems at once we need a rule for combining their state spaces.

\subsection{The tensor product}

\begin{postulate}[Composite systems]\label{post:composite}
  The state space of a composite physical system is the \emph{tensor product} of the state spaces of its components. If systems $1,\dots,n$ are prepared in the states $\ket{\psi_i}$, the joint state of the total system is
  $\ket{\psi_1}\tp\ket{\psi_2}\tp\cdots\tp\ket{\psi_n}$.
\end{postulate}

Recall the construction. If $\{\ket{e_i}\}$ and $\{\ket{f_j}\}$ are orthonormal bases of $\Hil_1$ and $\Hil_2$, then $\{\ket{e_i}\tp\ket{f_j}\}$ is an orthonormal basis of $\Hil_1\tp\Hil_2$, so
\[
  \dim(\Hil_1\tp\Hil_2)=\dim\Hil_1\cdot\dim\Hil_2,\qquad \braket{a\tp b}{c\tp d}=\braket ac\,\braket bd .
\]
One can simply accept the postulate as a basic law. A heuristic: the superposition principle says that if $\ket x,\ket y$ are states, so is $\alpha\ket x+\beta\ket y$. For composite systems, if $\ket A$ is a state of $A$ and $\ket B$ one of $B$,
there should be a joint state $\ket A\ket B$; imposing superposition on such product states gives precisely $\Hil_A\tp\Hil_B$. (This is not a derivation, but it shows the flavour.)

We write $\ket a\tp\ket b=\ket a\ket b=\ket{ab}$. For two qubits the computational basis is $\ket{00},\ket{01},\ket{10},\ket{11}$, and the general state is
\[
  \ket\psi=\sum_{i,j\in\{0,1\}}c_{ij}\ket{ij},\qquad\sum_{ij}\abs{c_{ij}}^2=1 .
\]
For operators, $(A\tp B)(\ket a\tp\ket b)=A\ket a\tp B\ket b$, and the matrix of $A\tp B$ is the \term{Kronecker product}, with $B$ varying fastest:
\[
  \begin{pmatrix}a&b\\c&d\end{pmatrix}\tp B=\begin{pmatrix}aB&bB\\cB&dB\end{pmatrix}.
\]
The rules one uses constantly are, for all operators of compatible size,
\[
  (A\tp B)(C\tp D)=AC\tp BD,\quad (A\tp B)^\dagger=A^\dagger\tp B^\dagger,\quad \Tr(A\tp B)=\Tr A\,\Tr B,
\]
and the tensor product of unitaries (resp.\ Hermitian operators, projectors) is unitary (resp.\ Hermitian, a projector). If $A\ket a=\lambda\ket a$ and $B\ket b=\mu\ket b$ then $(A\tp B)\ket{ab}=\lambda\mu\ket{ab}$:
the spectrum of $A\tp B$ consists of the products of eigenvalues. In a system of several qubits we use subscripts for the qubit acted on: in three qubits $X_2=\Id\tp X\tp\Id$, and in two qubits $X_1Z_2=X\tp Z$.

\begin{example}\label{ex:kron}
  Write $X\tp Z$ as a $4\times4$ matrix; compute $(H\tp H)\ket{00}$ and $(H\tp H)\ket{11}$.
\end{example}
\begin{solution}
  By the block rule, with $X=\begin{pmatrix}0&1\\1&0\end{pmatrix}$ and $B=Z$,
  \[
    X\tp Z=\begin{pmatrix}0\cdot Z&1\cdot Z\\1\cdot Z&0\cdot Z\end{pmatrix}=\begin{pmatrix}0&0&1&0\\0&0&0&-1\\1&0&0&0\\0&-1&0&0\end{pmatrix}.
  \]
  Next $(H\tp H)\ket{00}=H\ket0\tp H\ket0=\ket+\ket+=\tfrac12(\ket{00}+\ket{01}+\ket{10}+\ket{11})$, while
  $(H\tp H)\ket{11}=\ket-\ket-=\tfrac12(\ket{00}-\ket{01}-\ket{10}+\ket{11})$. In general $H\tp H\ket{xy}=\tfrac12\sum_{u,v}(-1)^{xu+yv}\ket{uv}$,
  and for $n$ qubits $H^{\tp n}\ket x=2^{-n/2}\sum_y(-1)^{x\cdot y}\ket y$: the Walsh--Hadamard transform, a first step of many algorithms.
\end{solution}

\paragraph{Exponential growth.}
$n$ qubits live in $\C^{2^n}$: a general state needs $2^n$ amplitudes. We computed in Example~\ref{ex:nqubits} what this means in memory for $n=50$.

\subsection{Entanglement}

A product state $\ket a\tp\ket b$ has a definite state for each part. But the postulate also allows superpositions of products, and some of these cannot be written as a single product.

\begin{definition}[Entanglement]\label{def:entangled}
  A state of a composite system that cannot be written as a product $\ket a\tp\ket b$ of states of its components is called \term{entangled}.
\end{definition}

The prototype is the two-qubit state
\[
  \ket{\Phi^+}=\frac{\ket{00}+\ket{11}}{\sqrt2},
\]
which has no single-qubit states $\ket a,\ket b$ with $\ket{\Phi^+}=\ket a\ket b$ (Theorem~\ref{thm:entangled-criterion} below). It is one of four \term{Bell states},
\[
  \bell xy=\frac{\ket{0y}+(-1)^x\ket{1\bar y}}{\sqrt2},\qquad x,y\in\{0,1\},\ \bar y=1-y,
\]
that is, $\bell00=\ket{\Phi^+}$, $\bell10=\frac{\ket{00}-\ket{11}}{\sqrt2}$, $\bell01=\frac{\ket{01}+\ket{10}}{\sqrt2}$ and $\bell11=\frac{\ket{01}-\ket{10}}{\sqrt2}$ (the last is the \emph{singlet}). They form an orthonormal basis of the two-qubit space (Exercise~\ref{ex:bell-basis}).

For two qubits there is a simple test. Arrange the amplitudes of $\ket\psi=\sum_{ij}c_{ij}\ket{ij}$ in a $2\times2$ matrix $C=(c_{ij})$.

\begin{theorem}[Criterion for two qubits]\label{thm:entangled-criterion}
  $\ket\psi$ is a product state if and only if $\rank C=1$, that is, if and only if $\det C=0$.
\end{theorem}
\begin{proof}
  If $\ket\psi=\ket a\ket b$ with $\ket a=\sum a_i\ket i$ and $\ket b=\sum b_j\ket j$, then $c_{ij}=a_ib_j$, so $C=ab^t$ has rank $1$ (it has rank $\le1$ and is nonzero). Conversely, a $2\times2$ matrix of rank $1$ has the form $ab^t$ for column vectors $a,b\neq0$, and then $\ket\psi=\ket a\ket b$ up to normalisation. For a $2\times2$ matrix, rank $1$ means $\det C=0$ (and $C\neq0$).
\end{proof}

For $\ket{\Phi^+}$ we get $C=\tfrac1{\sqrt2}\Id$ with $\det C=\tfrac12\neq0$: it is entangled. (The number $2\abs{\det C}$ is called the \emph{concurrence}; it equals $1$ for a Bell state and $0$ for a product state.)

\begin{example}\label{ex:entangled-or-not}
  Decide which are entangled: (a) $\ket+\ket-$; (b) $\frac12(\ket{00}+\ket{01}+\ket{10}+\ket{11})$; (c) $\frac1{\sqrt3}(\ket{00}+\ket{01}+\ket{10})$; (d) $\bell11$.
\end{example}
\begin{solution}
  (a) A product by construction, $\det C=0$. (b) $C=\tfrac12\begin{pmatrix}1&1\\1&1\end{pmatrix}$ has $\det C=0$: it is the product $\ket+\ket+$. (c) $C=\tfrac1{\sqrt3}\begin{pmatrix}1&1\\1&0\end{pmatrix}$ has $\det C=-\tfrac13\neq0$: entangled.
  (d) $C=\tfrac1{\sqrt2}\begin{pmatrix}0&1\\-1&0\end{pmatrix}$ has $\det C=\tfrac12\neq0$: entangled.
\end{solution}

\subsection{Measuring one part}

How does Postulate~3 act on a composite system? To measure only the first qubit with $\{M_m\}$, use the operators $M_m\tp\Id$. For a computational-basis measurement, $M_m=\ket m\!\bra m$.

\begin{example}\label{ex:partial}
  Measure the first qubit of $\ket{\Phi^+}$ in the computational basis. What happens to the second?
\end{example}
\begin{solution}
  With $M_m=\ket m\!\bra m\tp\Id$: $p(m)=\bra{\Phi^+}M_m\ket{\Phi^+}=\tfrac12$ for $m=0,1$, and the post-measurement state is $M_m\ket{\Phi^+}/\sqrt{p(m)}=\ket{mm}$.
  Each outcome by itself is a fair coin, but afterwards the second qubit is in the same state as the first: the outcomes are perfectly correlated. Moreover, because $\ket{\Phi^+}=\frac1{\sqrt2}(\ket{++}+\ket{--})$ (Exercise~\ref{ex:phiplus-pm}), measuring both qubits in the $\ket\pm$ basis also gives equal results.
  The correlation is not about $\ket0,\ket1$ in particular. Compare with a product state such as $\ket+\ket+$: measuring the first qubit tells nothing about the second.
\end{solution}

\begin{example}\label{ex:partial2}
  Measure the first qubit of $\ket\psi=\frac1{\sqrt3}(\ket{00}+\ket{01}+\ket{10})$ in the computational basis.
\end{example}
\begin{solution}
  Outcome $0$: $M_0\ket\psi=\frac1{\sqrt3}(\ket{00}+\ket{01})$, so $p(0)=\tfrac23$ and the state afterwards is $\frac1{\sqrt2}(\ket{00}+\ket{01})=\ket0\ket+$.
  Outcome $1$: $M_1\ket\psi=\frac1{\sqrt3}\ket{10}$, so $p(1)=\tfrac13$ and the state is $\ket{10}$. The outcome of the first qubit changes the state of the second, which is what entanglement (Example~\ref{ex:entangled-or-not}(c)) means operationally.
\end{solution}

\subsection{General measurements from projective ones}

Using composite systems we can show that general measurements need nothing beyond unitaries and projective measurements.

\begin{proposition}\label{prop:ancilla}
  Every measurement $\{M_m\}$ on a system $Q$ can be realised by a unitary on $Q$ together with an ancilla, followed by a projective measurement of the ancilla.
\end{proposition}
\begin{proof}
  Let $M$ be an auxiliary system with orthonormal basis $\{\ket m\}$, one vector per outcome, prepared in a fixed state $\ket0_M$. Define, on states of the form $\ket\psi\ket0_M$,
  \[
    U\bigl(\ket\psi\ket0_M\bigr)=\sum_mM_m\ket\psi\ket m .
  \]
  This preserves inner products: $\bra\varphi\bra0U^\dagger U\ket\psi\ket0=\sum_{m,m'}\bra\varphi M_m^\dagger M_{m'}\ket\psi\braket m{m'}=\sum_m\bra\varphi M_m^\dagger M_m\ket\psi=\braket\varphi\psi$ by completeness.
  A map defined on a subspace that preserves inner products extends to a unitary on the whole space (complete orthonormal bases of the subspace and of its image, and match the orthogonal complements, which have the same dimension). Now measure the ancilla projectively with $P_m=\Id_Q\tp\ket m\!\bra m$. The probability of outcome $m$ is
  $\norm{(\Id\tp\ket m\!\bra m)\sum_{m'}M_{m'}\ket\psi\ket{m'}}^2=\bra\psi M_m^\dagger M_m\ket\psi$, and the joint state afterwards is $M_m\ket\psi\ket m/\sqrt{p(m)}$: $Q$ ends in exactly the state prescribed by Postulate~3.
\end{proof}

So general measurements are projective measurements on a larger system: a nice use of Postulate~4. This is why quantum circuits may restrict themselves to unitaries, ancilla qubits and measurements in the computational basis.

\subsection{The postulates in one glance}

\begin{center}
\small
\begin{tabular}{@{}llll@{}}
  \toprule
  Postulate & Object & Key formula & Role\\
  \midrule
  1. State & unit vector $\ket\psi\in\Hil$ & $\braket\psi\psi=1$ & prepare\\
  2. Evolution & unitary $U$; Hamiltonian $H$ & $\ket{\psi'}=U\ket\psi$; $\ii\hbar\tfrac{d}{dt}\ket\psi=H\ket\psi$ & rotate (reversible)\\
  3. Measurement & $\{M_m\}$; observables; POVMs & $p(m)=\bra\psi M_m^\dagger M_m\ket\psi$ & read (irreversible)\\
  4. Composition & $\Hil_1\tp\Hil_2$ & $\ket{\psi_1}\tp\ket{\psi_2}$, entanglement & combine systems\\
  \bottomrule
\end{tabular}
\end{center}

The strangeness of the framework in one paragraph: we can never observe the state vector directly. It is like a game of chess in which you may only learn the rank of each piece, never the piece. In classical physics, energy,
position and velocity are directly accessible and observing does not change the system; in quantum mechanics the fundamental object is a hidden vector, and observation is invasive. Can quantum mechanics be reformulated in a more classical way,
with hidden but definite properties? Section~\ref{sec:bell} shows that every such attempt is doomed. The good news is that this counter-intuitive behaviour is exactly what quantum computation exploits.

\begin{exerciseset}
  \exercise{1} Compute $\ket+\tp\ket-$ as a column vector, $(H\tp H)\ket{00}$ and $(H\tp H)\ket{11}$, and find the eigenvalues of $X\tp Z$.

  \exercise{2} Prove the rules $(A\tp B)(C\tp D)=AC\tp BD$, $(A\tp B)^\dagger=A^\dagger\tp B^\dagger$ and $\Tr(A\tp B)=\Tr A\Tr B$. Deduce that the tensor product of unitaries is unitary and of Hermitian operators is Hermitian.

  \exercise{2} Show that $H\tp H\ket{xy}=\frac12\sum_{u,v}(-1)^{xu+yv}\ket{uv}$, and generalise to $n$ qubits (the Walsh--Hadamard transform).

  \exercise{1} Use the determinant criterion to classify: $\ket+\ket-$; $\frac12(\ket{00}+\ket{01}+\ket{10}+\ket{11})$; $\frac12(\ket{00}+\ket{01}+\ket{10}-\ket{11})$; the four Bell states. Prove directly that $\ket{\Phi^+}\neq\ket a\ket b$ for all single-qubit states $\ket a,\ket b$.

  \exercise{1} Show that the four Bell states form an orthonormal basis of the two-qubit state space.\label{ex:bell-basis}

  \exercise{1} Show $\ket{\Phi^+}=\frac1{\sqrt2}(\ket{++}+\ket{--})$ and write $\bell11$ in the basis $\ket{\pm}\ket{\pm}$.\label{ex:phiplus-pm}

  \exercise{2} Show that the average value of $X_1Z_2$ in $\ket{\Phi^+}$ is zero. Compute $\langle X\tp X\rangle,\langle Y\tp Y\rangle,\langle Z\tp Z\rangle$ in $\ket{\Phi^+}$ and in $\bell11$. What do you observe?

  \exercise{2} Let $\CNOT=\ket0\!\bra0\tp\Id+\ket1\!\bra1\tp X$. Write it as a $4\times4$ matrix and prove it is unitary and involutory. Compute $\CNOT(H\tp\Id)\ket{xy}$ for the four basis states. What do you find? (This is the standard circuit that prepares the Bell states.)

  \exercise{2} Let $\ket\psi=\frac1{\sqrt3}(\ket{00}+\ket{01}+\ket{10})$. Measure the first qubit in the computational basis with operators $\ket m\!\bra m\tp\Id$ and compute the probabilities and post-measurement states. Repeat with $\ket+\ket-$ and with $\bell01$.

  \exercise{2} Show that the amplitudes of an $n$-qubit product state are determined by $2n$ real numbers up to phases, whereas a general state of $n$ qubits needs $2^{n+1}-2$. What does the gap say about entanglement?

  \exercise{3} For $\ket\psi=\sum_{ij}c_{ij}\ket{ij}$ and unitaries $U,V$, show that $(U\tp V)\ket\psi$ has coefficient matrix $UCV^t$, and deduce that the rank of $C$ does not change under local unitaries.

  \exercise{3} Prove the claim in the proof of Proposition~\ref{prop:ancilla}: a linear map $U:W\to V$ ($W$ a subspace) with $\braket{Uw_1}{Uw_2}=\braket{w_1}{w_2}$ extends to a unitary $V\to V$. (Hint: complete orthonormal bases of $W$ and $U(W)$ with Gram--Schmidt.)

  \exercise{2} Take the measurement of Exercise~\ref{ex:epsilon-meas} with $\epsilon=\frac12$ and write the unitary $U$ with $U\ket\psi\ket0=M_0\ket\psi\ket0+M_1\ket\psi\ket1$ explicitly on the vectors $\ket{\psi}\ket0$, and complete it to a $4\times4$ unitary.
\end{exerciseset}

%======================================================================
\section{Application: superdense coding}\label{sec:superdense}
%======================================================================

\begin{objectives}
  You will see all four postulates working together in a short protocol: two classical bits are transmitted
  by sending a single qubit, using an entangled pair shared in advance.
\end{objectives}

\begin{recall}
  The Bell states and the controlled-NOT (Section~\ref{sec:composite}); the gates $X$, $Z$, $H$ (Section~\ref{sec:state-evolution}); the fact that orthonormal states can be told apart with certainty (Section~\ref{sec:measurement}).
\end{recall}

A qubit lives in a two-dimensional space, yet a single measurement of it yields only one classical bit. It is therefore surprising that one qubit can
\emph{carry} two bits. The catch is that Alice and Bob must share an entangled pair beforehand. The protocol, due to Bennett and Wiesner~\cite{bennett1992}, is the simplest illustration of entanglement as a resource, and it uses every postulate.

\paragraph{The protocol.}
Alice and Bob each hold one qubit of the state $\bell00=\frac1{\sqrt2}(\ket{00}+\ket{11})$, Alice the first and Bob the second (Postulates~1 and~4), prepared at some time in the past; afterwards they separate.
To send the two bits $xy\in\{00,01,10,11\}$ Alice applies to her qubit the gate $X^yZ^x$ (Postulate~2) and sends it to Bob. Bob now holds both qubits. He applies $\CNOT$ and then a Hadamard gate to the first qubit, and measures both qubits in the computational basis (Postulate~3). The outcome is $xy$.

\begin{center}
\small
\begin{tabular}{@{}ccc@{}}
  \toprule
  bits $xy$ & Alice applies & state of the pair\\
  \midrule
  $00$ & $\Id$ & $\bell00$\\
  $01$ & $X$ & $\bell01$ (up to a phase)\\
  $10$ & $Z$ & $\bell10$\\
  $11$ & $XZ$ & $\bell11$ (up to a phase)\\
  \bottomrule
\end{tabular}
\end{center}

\begin{figure}[htb]
  \centering
  \begin{tikzpicture}[>=stealth,
      box/.style={draw=navy,thick,fill=navy!8,minimum height=0.8cm,font=\small},
      cl/.style={draw=figacc,thick,fill=figacc!10,minimum height=0.8cm,font=\small}]
    \node[draw=ink,thick,fill=ink!8,minimum width=2.3cm,minimum height=0.8cm,font=\small] (src) at (0,0) {$\bell00$};
    \node[font=\footnotesize,ink] at (0,0.75) {shared entangled pair};
    \node[box,minimum width=1.6cm] (enc) at (4,1.3) {$X^yZ^x$};
    \node[font=\footnotesize,ink] at (4,1.95) {Alice};
    \node[box,minimum width=2.6cm] (dec) at (8,0) {$\CNOT$, then $H\tp\Id$};
    \node[font=\footnotesize,ink] at (8,-0.75) {Bob};
    \node[cl,minimum width=1.6cm] (m) at (11.2,0) {measure};
    \draw[->,thick] (1.15,0.25) -- (2.3,0.25) -- (2.3,1.3) -- (enc.west);
    \draw[->,thick] (enc.east) -- (5.9,1.3) -- (5.9,0.25) -- (6.7,0.25);
    \node[font=\footnotesize] at (6.55,1.55) {one qubit sent};
    \draw[->,thick] (1.15,-0.25) -- node[below,font=\footnotesize] {Bob's qubit, kept} (6.7,-0.25);
    \draw[->] (dec) -- (m);
    \draw[->,figacc] (m.north) -- ++(0,0.5) node[above,font=\small] {$x,y$};
  \end{tikzpicture}
  \caption{Superdense coding. Alice encodes two classical bits $x,y$ with a gate on her half of a shared entangled pair, and sends this single qubit to Bob, who decodes both bits.}
  \label{fig:superdense}
\end{figure}

\begin{proposition}\label{prop:superdense}
  Alice's operation maps $\bell00$ to $\bell xy$ up to a global phase, and Bob's circuit sends $\bell xy$ to $\ket x\ket y$, so Bob recovers $xy$ with certainty.
\end{proposition}
\begin{proof}
  \begin{steps}
    \item \emph{Alice's $Z^x$.} Acting on the first qubit, $Z^x\tp\Id$ gives $\frac1{\sqrt2}(\ket{00}+(-1)^x\ket{11})$.
    \item \emph{Alice's $X^y$.} Then $X^y\tp\Id$ gives $\frac1{\sqrt2}(\ket{y0}+(-1)^x\ket{\bar y1})$.
    \item \emph{Compare with the Bell state.} For $y=0$ this is $\bell x0$. For $y=1$ it is $\frac1{\sqrt2}(\ket{10}+(-1)^x\ket{01})=(-1)^x\bell x1$, equal to $\bell x1$ up to a global phase.
    \item \emph{Bob's $\CNOT$.} It flips the second qubit when the first is $\ket1$: $\bell xy=\frac{1}{\sqrt2}(\ket{0y}+(-1)^x\ket{1\bar y})\mapsto\frac1{\sqrt2}(\ket0+(-1)^x\ket1)\ket y$, which is $\ket+\ket y$ for $x=0$ and $\ket-\ket y$ for $x=1$.
    \item \emph{Bob's Hadamard} on the first qubit maps $\ket\pm\mapsto\ket{0},\ket{1}$, leaving $\ket x\ket y$.
    \item \emph{Measure} in the computational basis: $(x,y)$ is returned with probability $1$. \qedhere
  \end{steps}
\end{proof}

The bits $x,y$ are therefore the two values Bob reads from the final measurement. Equivalently, Bob has performed a measurement in the Bell basis, which can only be told apart, perfectly, because the four Bell states are orthogonal (Proposition~\ref{prop:distinguish-orthogonal}).

\begin{example}\label{ex:superdense-run}
  Alice wants to send $xy=10$. Follow the protocol.
\end{example}
\begin{solution}
  \begin{steps}
    \item \emph{Alice.} She applies $Z$ to her qubit: $\frac1{\sqrt2}(\ket{00}+\ket{11})\mapsto\frac1{\sqrt2}(\ket{00}-\ket{11})=\bell10$.
    \item \emph{Bob's $\CNOT$.} $\ket{00}\mapsto\ket{00}$ and $\ket{11}\mapsto\ket{10}$, giving $\frac1{\sqrt2}(\ket0-\ket1)\ket0=\ket-\ket0$.
    \item \emph{Bob's Hadamard.} It sends $\ket-\mapsto\ket1$ on the first qubit, so the state is $\ket{10}$.
    \item \emph{Result.} Bob reads $x=1$, $y=0$.
  \end{steps}
\end{solution}

\begin{remark}
  Could an eavesdropper, Eve, learn something by intercepting the qubit on its way to Bob? Not without Bob's half: by itself, Alice's qubit is in the maximally mixed state $\Id/2$ for all four messages,
  so it carries no information at all about $xy$ (Example~\ref{ex:reduced-bell} below). The information is stored in the \emph{correlations} between the two qubits, not in either one.
  Nor does the protocol beat Nature's accounting: it transmits two bits using one qubit \emph{sent}, but one entangled pair \emph{already shared}.
\end{remark}

\begin{lookahead}
  Superdense coding has a mirror image, \emph{teleportation} (Section~\ref{sec:teleportation}), in which a shared Bell pair and two classical bits let Alice transmit an unknown qubit to Bob. Both rest on the same fact: the four Bell states are perfectly distinguishable. In information-theoretic terms (Section~\ref{sec:vonneumann}), the shared pair carries a quantum mutual information of $2$ bits between Alice's and Bob's halves, which is exactly the number of classical bits that Alice can add to a single qubit sent.
\end{lookahead}

\begin{exerciseset}
  \exercise{1} Verify the table of operations: compute $\bigl(X^yZ^x\tp\Id\bigr)\bell00$ for each $xy$ and compare with $\bell xy$.

  \exercise{1} Run the protocol for $xy=01$ and $xy=11$, writing the state after each of Bob's gates.

  \exercise{2} Suppose Alice applies the Hadamard gate $H$ instead of one of the four gates. What does Bob measure? Which state does the pair end up in?

  \exercise{2} Let $E$ be any positive operator on Alice's qubit. Show that $\bra\psi(E\tp\Id)\ket\psi$ takes the same value for all four Bell states $\ket\psi$. (Hint: $\bra{\beta_{xy}}E\tp\Id\ket{\beta_{xy}}=\tfrac12\Tr E$.) What does this tell you about Eve?

  \exercise{3} Suppose Bob measures only the second qubit in the computational basis, before receiving Alice's. What can he learn about $xy$? Compare with what he learns after the full decoding.
\end{exerciseset}

%======================================================================
\section{The density operator}\label{sec:density}
%======================================================================

\begin{objectives}
  You will describe ensembles of states and subsystems of entangled systems with density operators; recognise pure and mixed states and place qubit states in the Bloch ball; see a state as a positive, normalised linear functional on observables (the Riesz theorem for operators); and compute reduced density operators with the partial trace.
\end{objectives}

\begin{recall}
  Positive operators (Section~\ref{sec:measurement}); the trace and the Hilbert--Schmidt inner product $\Tr(A^\dagger B)$ (Sections~\ref{sec:vector-spaces} and~\ref{sec:hilbert}); the Riesz theorem (Section~\ref{sec:state-evolution}); tensor products and measurements on one part of a system (Section~\ref{sec:composite}).
\end{recall}

So far a system has been described by a state vector, and that is enough if we know exactly how it was prepared. Two situations break this assumption. In the first, the preparation is itself random: a source emits $\ket{\psi_i}$ with probability $p_i$ and we do not know which.
In the second, the system is one \emph{part} of an entangled pair: we shall see that it does not have a state vector of its own. In both cases the right tool is the density operator.

\subsection{Ensembles of states}

\begin{definition}[Ensemble and density operator]\label{def:density}
  An \term{ensemble} $\{p_i,\ket{\psi_i}\}$ is a collection of states $\ket{\psi_i}$ with probabilities $p_i\ge0$, $\sum_ip_i=1$. Its \term{density operator} is
  \[
    \rho=\sum_ip_i\ket{\psi_i}\!\bra{\psi_i}.
  \]
  If a single $p_i=1$ the state is \term{pure}, $\rho=\ket\psi\!\bra\psi$; otherwise it is \term{mixed}.
\end{definition}

Why is this the right object? Suppose the system is in $\ket{\psi_i}$ with probability $p_i$ and we measure $\{M_m\}$.
\begin{steps}
  \item \emph{Probabilities.} By total probability, $p(m)=\sum_ip_i\,p(m\mid i)=\sum_ip_i\bra{\psi_i}M_m^\dagger M_m\ket{\psi_i}=\Tr\bigl(M_m^\dagger M_m\,\rho\bigr)$, using $\bra\psi A\ket\psi=\Tr(A\ket\psi\!\bra\psi)$.
  \item \emph{After the measurement.} If the state was $\ket{\psi_i}$ it becomes $M_m\ket{\psi_i}$ (normalised), so the post-measurement ensemble has density operator $M_m\rho M_m^\dagger/p(m)$.
  \item \emph{Evolution.} Under a unitary each $\ket{\psi_i}\mapsto U\ket{\psi_i}$, so $\rho\mapsto U\rho U^\dagger$.
\end{steps}
We have shown the following translation of the postulates.

\begin{proposition}[Postulates for density operators]\label{prop:postulates-rho}
  In terms of $\rho$: (2) a unitary evolution takes $\rho\mapsto U\rho U^\dagger$; (3) a measurement $\{M_m\}$ has outcome probabilities $p(m)=\Tr(M_m^\dagger M_m\rho)$ and post-measurement states $\rho_m=M_m\rho M_m^\dagger/p(m)$; and the expected value of an observable $M$ is $\Tr(M\rho)$.
\end{proposition}

The crucial point is that \emph{the density operator contains everything that can be observed}. Different ensembles can have the same $\rho$, and then they are physically indistinguishable.

\begin{example}\label{ex:mixed-vs-superposition}
  Compare the ensemble $\bigl\{\tfrac12,\ket0;\ \tfrac12,\ket1\bigr\}$, the ensemble $\bigl\{\tfrac12,\ket+;\ \tfrac12,\ket-\bigr\}$, and the pure state $\ket+$.
\end{example}
\begin{solution}
  \begin{steps}
    \item \emph{First ensemble.} $\rho=\tfrac12\ket0\!\bra0+\tfrac12\ket1\!\bra1=\tfrac12\Id$.
    \item \emph{Second ensemble.} $\rho=\tfrac12\ket+\!\bra++\tfrac12\ket-\!\bra-=\tfrac12\Id$ as well, by the completeness relation. So the two ensembles are physically identical: no measurement distinguishes them, in spite of the fact that their members look very different.
    \item \emph{The pure state.} $\ket+\!\bra+=\tfrac12\begin{pmatrix}1&1\\1&1\end{pmatrix}\neq\tfrac12\Id$. It can be told apart from the mixtures: measuring $X$ gives $+1$ with certainty on $\ket+$, but with probability $\tfrac12$ on $\tfrac12\Id$.
  \end{steps}
  The \emph{superposition} $\ket+$ is not the same as a \emph{coin toss} between $\ket0$ and $\ket1$.
\end{solution}

\begin{example}[A three-state ensemble]\label{ex:three-ensemble}
  A source emits $\ket0$ with probability $\tfrac12$, $\ket1$ with probability $\tfrac14$ and $\ket+$ with probability $\tfrac14$. Find $\rho$, and the probability that a measurement of $X$ gives $+1$.
\end{example}
\begin{solution}
  \begin{steps}
    \item \emph{The density matrix.} $\rho=\tfrac12\begin{pmatrix}1&0\\0&0\end{pmatrix}+\tfrac14\begin{pmatrix}0&0\\0&1\end{pmatrix}+\tfrac14\cdot\tfrac12\begin{pmatrix}1&1\\1&1\end{pmatrix}=\tfrac18\begin{pmatrix}5&1\\1&3\end{pmatrix}$.
    \item \emph{Trace check.} $\Tr\rho=\tfrac18(5+3)=1$.
    \item \emph{The probability.} $P(X=+1)=\Tr(\ket+\!\bra+\rho)=\bra+\rho\ket+=\tfrac12\cdot\tfrac18(5+1+1+3)=\tfrac58$.
    \item \emph{Check by total probability.} $\tfrac12\cdot\tfrac12+\tfrac14\cdot\tfrac12+\tfrac14\cdot1=\tfrac58$.
  \end{steps}
\end{solution}

\subsection{Characterising density operators}

We can now forget ensembles and characterise density operators intrinsically.

\begin{theorem}[Characterisation]\label{thm:density-char}
  An operator $\rho$ is the density operator of some ensemble if and only if (a) $\Tr\rho=1$, and (b) $\rho$ is positive (Definition~\ref{def:positive}).
\end{theorem}
\begin{proof}
  \begin{steps}
    \item \emph{Necessity.} If $\rho=\sum_ip_i\ket{\psi_i}\!\bra{\psi_i}$, then $\Tr\rho=\sum_ip_i=1$ and $\bra\varphi\rho\ket\varphi=\sum_ip_i\abs{\braket\varphi{\psi_i}}^2\ge0$.
    \item \emph{Sufficiency.} Suppose (a) and (b). By Theorem~\ref{thm:positive}, $\rho$ is Hermitian with eigenvalues $\lambda_j\ge0$: $\rho=\sum_j\lambda_j\ket j\!\bra j$ for an orthonormal basis $\ket j$.
    \item \emph{The ensemble.} $\sum_j\lambda_j=\Tr\rho=1$, so $\{\lambda_j,\ket j\}$ is an ensemble, and its density operator is $\rho$. \qedhere
  \end{steps}
\end{proof}

When is a state pure? Since $\Tr\rho^2=\sum_j\lambda_j^2\le\bigl(\sum_j\lambda_j\bigr)^2=1$, with equality iff there is a single non-zero eigenvalue, we get the \emph{purity test}:
\[
  \Tr\rho^2\le1,\quad\text{with equality if and only if }\rho\text{ is pure.}
\]

\paragraph{The Bloch ball.}
Any $2\times2$ Hermitian matrix of trace $1$ is $\rho=\tfrac12(\Id+\vec r\cdot\vec\sigma)$ with $\vec r\in\R^3$ (Theorem~\ref{thm:paulibasis}), and its eigenvalues are $\tfrac12(1\pm\abs{\vec r})$. Positivity therefore means $\abs{\vec r}\le1$: \emph{qubit density operators are exactly the points of the unit ball},
with pure states on the boundary sphere (Proposition~\ref{prop:bloch}) and the maximally mixed state $\tfrac12\Id$ at the centre (Figure~\ref{fig:blochball}). Mixing states averages their Bloch vectors, and $\Tr\rho^2=\tfrac12(1+\abs{\vec r}^2)$.

\begin{figure}[htb]
  \centering
  \begin{tikzpicture}[>=stealth,scale=1.3]
    \draw[ink,thick,fill=navy!5] (0,0) circle (2);
    \draw[ink!50,dashed] (-2,0) arc (180:0:2 and 0.6);
    \draw[ink!60] (-2,0) arc (180:360:2 and 0.6);
    \draw[->,ink!70] (0,0) -- (0,2.5) node[above] {$z$};
    \fill[navy] (0,2) circle (2pt) node[above left,font=\small] {$\ket0$};
    \fill[navy] (0,-2) circle (2pt) node[below left,font=\small] {$\ket1$};
    \fill[ink] (0,0) circle (2pt) node[below left,font=\small] {$\tfrac12\Id$};
    \draw[very thick,figacc,->] (0,0) -- (0.9,0.9);
    \fill[figacc] (0.9,0.9) circle (2pt) node[right,font=\small] {mixed: $\abs{\vec r}<1$};
    \fill[navy] (1.28,1.5) circle (2pt) node[above right,font=\small] {pure: $\abs{\vec r}=1$};
  \end{tikzpicture}
  \caption{The Bloch ball. Pure states are on the surface; mixed states are inside, with the maximally mixed state $\Id/2$ at the centre.}
  \label{fig:blochball}
\end{figure}

\begin{example}\label{ex:mixture-bloch}
  Find the density matrix and Bloch vector of the ensemble $\bigl\{\tfrac12,\ket0;\ \tfrac12,\ket+\bigr\}$. Is it pure?
\end{example}
\begin{solution}
  \begin{steps}
    \item \emph{The matrix.} $\rho=\tfrac12\ket0\!\bra0+\tfrac12\ket+\!\bra+=\tfrac12\begin{pmatrix}1&0\\0&0\end{pmatrix}+\tfrac14\begin{pmatrix}1&1\\1&1\end{pmatrix}=\tfrac14\begin{pmatrix}3&1\\1&1\end{pmatrix}$.
    \item \emph{The Bloch vector.} Averaging $(0,0,1)$ and $(1,0,0)$ gives $\vec r=(\tfrac12,0,\tfrac12)$, with $\abs{\vec r}=\tfrac1{\sqrt2}<1$: mixed.
    \item \emph{Eigenvalues.} $\tfrac12(1\pm\tfrac1{\sqrt2})\approx0.85,\ 0.15$.
    \item \emph{Purity check.} $\Tr\rho^2=\tfrac1{16}(9+1+1+1)=\tfrac34=\tfrac12(1+\tfrac12)$.
  \end{steps}
\end{solution}

\begin{remark}
  A given $\rho$ arises from many different ensembles, as Example~\ref{ex:mixed-vs-superposition} showed. It can be shown that two ensembles $\{p_i,\ket{\psi_i}\}$ and $\{q_j,\ket{\varphi_j}\}$ give the same $\rho$ exactly when $\sqrt{p_i}\ket{\psi_i}=\sum_ju_{ij}\sqrt{q_j}\ket{\varphi_j}$ for a unitary matrix $(u_{ij})$ (padding the shorter ensemble with zero vectors). This is the \emph{unitary freedom} of ensembles.
\end{remark}

\subsection{States as functionals: the Riesz theorem for operators}

There is a second, more abstract, way to see what a state is, and it uses the Riesz theorem of Section~\ref{sec:state-evolution} on a space different from $\Hil$. What an experimentalist ultimately has is a rule that assigns to every observable $A$ its expectation value $\omega(A)$. That rule is a \emph{linear functional on the space of operators}, and it is positive and normalised:
\begin{steps}
  \item \emph{linear:} $\omega(aA+bB)=a\,\omega(A)+b\,\omega(B)$;
  \item \emph{positive:} $\omega(A)\ge0$ whenever $A\ge0$ (an event cannot have negative expected count);
  \item \emph{normalised:} $\omega(\Id)=1$.
\end{steps}

\begin{theorem}[Riesz theorem for operators]\label{thm:riesz-operators}
  Let $\omega:L(\Hil)\to\C$ be a linear functional. Then there is a unique operator $\rho$ such that $\omega(A)=\Tr(\rho A)$ for all $A$. Moreover, $\omega$ is positive and normalised if and only if $\rho$ is a density operator.
\end{theorem}
\begin{proof}
  \begin{steps}
    \item \emph{A Hilbert space of operators.} $L(\Hil)$ with the Hilbert--Schmidt inner product $\langle T,A\rangle=\Tr(T^\dagger A)$ is a finite-dimensional Hilbert space (Exercise~\ref{ex:hs-inner}).
    \item \emph{Apply Riesz.} By Theorem~\ref{thm:riesz} there is a unique $T$ with $\omega(A)=\langle T,A\rangle=\Tr(T^\dagger A)$. Put $\rho=T^\dagger$. This proves existence and uniqueness.
    \item \emph{Normalisation.} $\omega(\Id)=\Tr\rho$, so $\omega(\Id)=1$ iff $\Tr\rho=1$.
    \item \emph{Positivity of $\rho$.} If $\omega$ is positive, then for every $\ket\psi$ the operator $\ket\psi\!\bra\psi$ is positive, so $0\le\omega(\ket\psi\!\bra\psi)=\Tr(\rho\ket\psi\!\bra\psi)=\bra\psi\rho\ket\psi$. Thus $\rho\ge0$.
    \item \emph{Conversely,} if $\rho\ge0$ and $A\ge0$, then $\Tr(\rho A)=\Tr\bigl(\sqrt\rho A\sqrt\rho\bigr)\ge0$ because $\sqrt\rho A\sqrt\rho\ge0$ (Exercise~\ref{ex:conjugation-positive}) and positive operators have non-negative trace.
    \item \emph{Conclusion.} By Theorem~\ref{thm:density-char}, $\omega$ is positive and normalised iff $\rho$ is a density operator. \qedhere
  \end{steps}
\end{proof}

So ``a state'' has three equivalent descriptions: a unit vector up to phase (pure states), a density operator, and a positive normalised functional $A\mapsto\Tr(\rho A)$ on observables. The last is the starting point of the algebraic approach to quantum theory, in which states are \emph{defined} as such functionals.

\begin{example}\label{ex:riesz-state}
  Find the density operator representing the functional $\omega(A)=\tfrac12\bra0A\ket0+\tfrac12\bra+A\ket+$ and check that it is a state.
\end{example}
\begin{solution}
  \begin{steps}
    \item \emph{Rewrite with the trace.} $\bra\psi A\ket\psi=\Tr(\ket\psi\!\bra\psi A)$, so $\omega(A)=\Tr\bigl(\tfrac12(\ket0\!\bra0+\ket+\!\bra+)A\bigr)$.
    \item \emph{The representer.} $\rho=\tfrac12\ket0\!\bra0+\tfrac12\ket+\!\bra+=\tfrac14\begin{pmatrix}3&1\\1&1\end{pmatrix}$, the state of Example~\ref{ex:mixture-bloch}.
    \item \emph{Check.} $\Tr\rho=1$ and $\rho\ge0$ (eigenvalues $0.85,0.15$), so $\omega$ is positive and normalised.
  \end{steps}
\end{solution}

\begin{lookahead}
  The converse question, ``which probability assignments to measurement outcomes are consistent?'', has a remarkable answer, Gleason's theorem: in dimension $\ge3$, any assignment of probabilities to projectors that is additive on orthogonal projectors must be of the form $P\mapsto\Tr(\rho P)$. Probabilities \emph{are} density operators, and the Born rule is not an extra assumption.
\end{lookahead}

\subsection{The reduced density operator}

The deepest use of $\rho$ is describing a \emph{part} of a composite system. Let $\rho^{AB}$ be a density operator on $\Hil_A\tp\Hil_B$.

\begin{definition}[Partial trace and reduced density operator]\label{def:ptrace}
  The \term{partial trace} over $B$ is the linear map determined by
  \[
    \Tr_B\bigl(\ket{a_1}\!\bra{a_2}\tp\ket{b_1}\!\bra{b_2}\bigr)=\ket{a_1}\!\bra{a_2}\;\braket{b_2}{b_1}.
  \]
  The \term{reduced density operator} of $A$ is $\rho^A=\Tr_B\rho^{AB}$.
\end{definition}

It is the right description of $A$ alone because it reproduces every statistic of measurements on $A$.

\begin{proposition}\label{prop:reduced}
  For every observable $M$ on $A$, $\ \Tr\bigl((M\tp\Id)\rho^{AB}\bigr)=\Tr\bigl(M\rho^A\bigr)$.
\end{proposition}
\begin{proof}
  \begin{steps}
    \item \emph{Reduce to simple tensors.} Both sides are linear in $\rho^{AB}$, so it is enough to check on $\rho^{AB}=\ket{a_1}\!\bra{a_2}\tp\ket{b_1}\!\bra{b_2}$.
    \item \emph{Left side.} $\Tr\bigl(M\ket{a_1}\!\bra{a_2}\bigr)\Tr\bigl(\ket{b_1}\!\bra{b_2}\bigr)=\bra{a_2}M\ket{a_1}\braket{b_2}{b_1}$.
    \item \emph{Right side.} $\Tr\bigl(M\ket{a_1}\!\bra{a_2}\bigr)\braket{b_2}{b_1}$, the same. \qedhere
  \end{steps}
\end{proof}

For a pure state $\ket\psi=\sum_{ij}c_{ij}\ket{ij}$ with coefficient matrix $C=(c_{ij})$, a direct computation gives the handy formula $\rho^A=CC^\dagger$ (and $\rho^B=(C^\dagger C)^t$).

\begin{example}\label{ex:reduced-bell}
  Find the reduced density operators of each qubit of every Bell state, and of $\ket a\ket b$.
\end{example}
\begin{solution}
  \begin{steps}
    \item \emph{The state.} For $\ket{\Phi^+}=\bell00$, $\rho^{AB}=\tfrac12\bigl(\ket{00}\!\bra{00}+\ket{00}\!\bra{11}+\ket{11}\!\bra{00}+\ket{11}\!\bra{11}\bigr)$.
    \item \emph{Partial trace over $B$.} The cross terms contain $\braket10=0$ or $\braket01=0$ and vanish, while $\braket00=\braket11=1$: $\rho^A=\tfrac12\bigl(\ket0\!\bra0+\ket1\!\bra1\bigr)=\tfrac12\Id$. Likewise $\rho^B=\tfrac12\Id$.
    \item \emph{The other Bell states.} Their coefficient matrices $C=U/\sqrt2$ are $1/\sqrt2$ times a unitary $U$, so $CC^\dagger=\tfrac12\Id$.
    \item \emph{A product state.} For $\ket a\ket b$: $\rho^A=\ket a\!\bra a\,\braket bb=\ket a\!\bra a$, pure.
  \end{steps}
\end{solution}

This is a remarkable fact. The two qubits together are in a \emph{pure} state, with no uncertainty whatsoever, yet each qubit by itself is maximally mixed, as random as it can be. All the information is in the correlations between the parts. The whole is more definite than its parts, and \emph{the state of a subsystem of an entangled system is mixed}.

\begin{example}\label{ex:reduced-3term}
  Compute $\rho^A$ for $\ket\psi=\frac1{\sqrt3}(\ket{00}+\ket{01}+\ket{10})$ and its eigenvalues.
\end{example}
\begin{solution}
  \begin{steps}
    \item \emph{Coefficient matrix.} $C=\frac1{\sqrt3}\begin{pmatrix}1&1\\1&0\end{pmatrix}$.
    \item \emph{Reduced state.} $\rho^A=CC^\dagger=\tfrac13\begin{pmatrix}2&1\\1&1\end{pmatrix}$, with $\Tr\rho^A=1$.
    \item \emph{Eigenvalues.} The characteristic polynomial is $\lambda^2-\lambda+\tfrac19$, so $\lambda=\tfrac12\bigl(1\pm\tfrac{\sqrt5}3\bigr)=\tfrac{3\pm\sqrt5}6\approx0.873,\ 0.127$.
    \item \emph{Conclusion.} Both are positive and below $1$: $A$ is mixed, so $\ket\psi$ is entangled, which agrees with Example~\ref{ex:entangled-or-not}(c). (These eigenvalues are the squares of the Schmidt coefficients of Section~\ref{sec:schmidt}, and, as we will see in Section~\ref{sec:vonneumann}, determine the entanglement entropy.)
  \end{steps}
\end{solution}

\begin{example}[Reduced state of a product of mixed states]\label{ex:reduced-product}
  Let $\rho^{AB}=\rho_1\otimes\rho_2$ with $\Tr\rho_1=\Tr\rho_2=1$. Find $\rho^A$ and $\rho^B$.
\end{example}
\begin{solution}
  \begin{steps}
    \item \emph{Use linearity.} Write $\rho_1=\sum_{a,a'}r_{aa'}\ket a\!\bra{a'}$ and $\rho_2=\sum_{b,b'}s_{bb'}\ket b\!\bra{b'}$, so $\rho_1\otimes\rho_2=\sum r_{aa'}s_{bb'}\ket a\!\bra{a'}\otimes\ket b\!\bra{b'}$.
    \item \emph{Trace out $B$.} Each term contributes $r_{aa'}s_{bb'}\ket a\!\bra{a'}\braket{b'}{b}=r_{aa'}s_{bb}\delta_{bb'}\ket a\!\bra{a'}$, and summing over $b$ gives $\Tr\rho_2\cdot\rho_1=\rho_1$.
    \item \emph{Result.} $\rho^A=\rho_1$ and, symmetrically, $\rho^B=\rho_2$: the parts of a product keep their own states.
  \end{steps}
\end{solution}

\begin{exerciseset}
  \exercise{1} Compute the density matrices of the ensembles $\bigl\{\tfrac12,\ket0;\tfrac12,\ket1\bigr\}$, $\bigl\{\tfrac12,\ket+;\tfrac12,\ket-\bigr\}$ and $\bigl\{\tfrac12,\ket{+\ii};\tfrac12,\ket{-\ii}\bigr\}$. Is there a measurement that tells them apart?

  \exercise{1} A source emits $\ket0$ with probability $\tfrac23$ and $\ket{-}$ with probability $\tfrac13$. Find $\rho$, its eigenvalues, and the probability of the outcome $+1$ for $Z$ and for $X$.

  \exercise{2} Show that the density operator of an ensemble transforms as $\rho\mapsto U\rho U^\dagger$ and, under a measurement, $\rho\mapsto M_m\rho M_m^\dagger/p(m)$ with $p(m)=\Tr(M_m^\dagger M_m\rho)$ (Proposition~\ref{prop:postulates-rho}).

  \exercise{2}\label{ex:positive-hermitian-sec} Show that a positive operator is Hermitian. (Hint: write $A=B+\ii C$ with $B,C$ Hermitian and use $\bra\psi A\ket\psi\in\R$.)

  \exercise{2} Prove the purity test: $\Tr\rho^2\le1$ with equality iff $\rho$ is pure. Show that for a qubit $\Tr\rho^2=\tfrac12(1+\abs{\vec r}^2)$.

  \exercise{2} Find the density matrix, Bloch vector and purity of $\bigl\{\tfrac13,\ket0;\ \tfrac23,\ket1\bigr\}$ and of $\bigl\{\tfrac12,\ket0;\ \tfrac12,\ket+\bigr\}$ (compare Example~\ref{ex:mixture-bloch}).

  \exercise{3} Show that a state $\rho=\tfrac12(\Id+\vec r\cdot\vec\sigma)$ evolves under $U=e^{-\ii\theta\vec n\cdot\vec\sigma/2}$ to the state whose Bloch vector is $\vec r$ rotated by $\theta$ about $\vec n$. (Hint: it suffices to check for $\vec n=\hat z$.)

  \exercise{2}\label{ex:hs-inner} Show that $\langle A,B\rangle=\Tr(A^\dagger B)$ is an inner product on $L(\Hil)$, and that $\{\tfrac1{\sqrt2}\sigma_\mu\}_{\mu=0}^3$ is an orthonormal basis of $L(\C^2)$. Use it to compute the representer, in the sense of Theorem~\ref{thm:riesz-operators}, of $\omega(A)=\tfrac12\Tr(A)$.

  \exercise{2}\label{ex:conjugation-positive} Show that if $A\ge0$ then $CAC^\dagger\ge0$ for every operator $C$, and that a positive operator has non-negative trace. Deduce that $\Tr(\rho A)\ge0$ whenever $\rho,A\ge0$.

  \exercise{2} Find the density operator representing the functional $\omega(A)=\tfrac13\bra0A\ket0+\tfrac23\bra1A\ket1$, and the one representing $\omega(A)=\bra{+\ii}A\ket{+\ii}$. Which are pure?

  \exercise{3} Show that a linear functional $\omega$ with $\omega(\Id)=1$ that is \emph{not} positive can give a ``negative expectation value''. Construct one: $\omega(A)=\Tr(\rho A)$ with $\rho=\diag(\tfrac32,-\tfrac12)$, and find a positive $A$ with $\omega(A)<0$.

  \exercise{2} Let $\ket\psi=\ket a\ket b$ be a product state. Show $\rho^A=\ket a\!\bra a$. Compute the reduced density operators of each qubit for $\ket{+}\ket{0}$, the four Bell states, and $\frac1{\sqrt2}(\ket{00}+\ket{01})$.

  \exercise{2} Prove $\rho^A=CC^\dagger$ for $\ket\psi=\sum c_{ij}\ket{ij}$, and show $\rho^B=(C^\dagger C)^t$. Deduce that $\rho^A$ is pure if and only if $\rank C=1$, that is, iff $\ket\psi$ is a product state.

  \exercise{3} Show that if $\ket{\psi}$ is a pure state of $AB$ then $\rho^A$ and $\rho^B$ have the same non-zero eigenvalues. (We will see why in Section~\ref{sec:schmidt}.)

  \exercise{2} Let $\rho^{AB}=\rho_1\otimes\rho_2$. Show that $\rho^A=\rho_1$ (Example~\ref{ex:reduced-product}) and compute $\rho^A,\rho^B$ for $\rho^{AB}=\tfrac12\ket{00}\!\bra{00}+\tfrac12\ket{11}\!\bra{11}$. Is this state a product?

  \exercise{3} Let Charlie prepare $\bell00$ and send one qubit to Alice and one to Bob. Alice applies any unitary $U$ to hers. Show that Bob's reduced density operator is unchanged, so that Alice cannot send Bob a message just by acting on her qubit.
\end{exerciseset}

%======================================================================
\section{Schmidt decomposition and purification}\label{sec:schmidt}
%======================================================================

\begin{objectives}
  You will write any pure state of two systems in the Schmidt form, read off whether it is entangled from its
  Schmidt coefficients, and show that every mixed state is a part of a pure state of a larger system.
\end{objectives}

The determinant test of Theorem~\ref{thm:entangled-criterion} detects entanglement for two qubits. For larger systems, and for a more refined answer to ``how entangled?'', we need a canonical form for a bipartite pure state. Its proof rests on one tool from matrix analysis,
which we use without proof (see for example Horn and Johnson~\cite{horn2013}).

\begin{theorem}[Singular value decomposition]\label{thm:svd}
  Every square complex matrix $C$ can be written $C=UDV^\dagger$ with $U,V$ unitary and $D$ diagonal with non-negative entries, the \term{singular values} of $C$. These are the square roots of the eigenvalues of $C^\dagger C$.
\end{theorem}

\subsection{The Schmidt decomposition}

\begin{theorem}[Schmidt decomposition]\label{thm:schmidt}
  Let $\ket\psi$ be a pure state of a composite system $AB$ with $\dim\Hil_A=\dim\Hil_B=d$. Then there exist orthonormal bases $\{\ket{i_A}\}$ of $\Hil_A$ and $\{\ket{i_B}\}$ of $\Hil_B$, and numbers $\lambda_i\ge0$ with $\sum_i\lambda_i^2=1$, such that
  \[
    \ket\psi=\sum_{i=1}^{d}\lambda_i\ket{i_A}\ket{i_B}.
  \]
\end{theorem}
\begin{proof}
  Write $\ket\psi=\sum_{jk}c_{jk}\ket j\ket k$ in any orthonormal bases and use the singular value decomposition $C=UDV^\dagger$, that is $c_{jk}=\sum_iU_{ji}\lambda_i\overline{V_{ki}}$. Then
  \[
    \ket\psi=\sum_i\lambda_i\Bigl(\sum_jU_{ji}\ket j\Bigr)\Bigl(\sum_k\overline{V_{ki}}\ket k\Bigr)=\sum_i\lambda_i\ket{i_A}\ket{i_B},
  \]
  where $\ket{i_A}=\sum_jU_{ji}\ket j$ are the columns of the unitary $U$ and $\ket{i_B}=\sum_k\overline{V_{ki}}\ket k$ the columns of the unitary $\bar V$, so both families are orthonormal. Finally $\sum_i\lambda_i^2=\braket\psi\psi=1$.
\end{proof}

(If $\dim\Hil_A\neq\dim\Hil_B$, the same holds with the number of terms equal to the smaller dimension.) The numbers $\lambda_i$ are the \term{Schmidt coefficients}, and the number of non-zero ones is the \term{Schmidt number}. The bases depend on the state, unlike in the general expansion $\sum c_{jk}\ket{jk}$, which needs a double sum.

\begin{corollary}\label{cor:schmidt}
  Let $\ket\psi=\sum_i\lambda_i\ket{i_A}\ket{i_B}$ be a Schmidt decomposition. Then
  \begin{enumerate}
    \item[(a)] $\rho^A=\sum_i\lambda_i^2\ket{i_A}\!\bra{i_A}$ and $\rho^B=\sum_i\lambda_i^2\ket{i_B}\!\bra{i_B}$: both reduced states have the same eigenvalues $\lambda_i^2$;
    \item[(b)] $\ket\psi$ is a product state if and only if its Schmidt number is $1$;
    \item[(c)] $\ket\psi$ is \term{maximally entangled} if all $d$ Schmidt coefficients equal $1/\sqrt d$, which happens iff $\rho^A=\Id/d$.
  \end{enumerate}
\end{corollary}
\begin{proof}
  (a) Take $\Tr_B$ of $\sum_{i,i'}\lambda_i\lambda_{i'}\ket{i_A}\!\bra{i'_A}\tp\ket{i_B}\!\bra{i'_B}$ and use $\braket{i'_B}{i_B}=\delta_{ii'}$. (b) A single term $\lambda_1=1$ is $\ket{1_A}\ket{1_B}$; conversely, a product state has $\rho^A$ pure, with one non-zero eigenvalue $1$. (c) is read off from (a).
\end{proof}

So \emph{the spectrum of the reduced state is a complete invariant of the entanglement of a pure state}: how mixed the parts are tells how entangled the whole is.

\begin{example}\label{ex:schmidt}
  Find the Schmidt coefficients of $\bell00$ and of $\ket\psi=\frac1{\sqrt3}(\ket{00}+\ket{01}+\ket{10})$.
\end{example}
\begin{solution}
  For $\bell00=\frac1{\sqrt2}(\ket{00}+\ket{11})$ the decomposition is already in Schmidt form: $\lambda_1=\lambda_2=\tfrac1{\sqrt2}$, with bases $\{\ket0,\ket1\}$ on both sides. It is maximally entangled.

  For $\ket\psi$, the reduced state $\rho^A=\tfrac13\begin{pmatrix}2&1\\1&1\end{pmatrix}$ has eigenvalues $\tfrac{3\pm\sqrt5}{6}$ (Example~\ref{ex:reduced-3term}), so by Corollary~\ref{cor:schmidt}(a)
  \[
    \lambda_{1,2}=\sqrt{\tfrac{3\pm\sqrt5}{6}}=\frac{\varphi}{\sqrt3},\ \frac1{\varphi\sqrt3}\ \approx\ 0.934,\ 0.357,
  \]
  where $\varphi=\frac{1+\sqrt5}{2}$ is the golden ratio (indeed $\varphi^2=\frac{3+\sqrt5}{2}$ and $\varphi^{-2}=\frac{3-\sqrt5}2$). \emph{Check:} $0.934^2+0.357^2=0.873+0.127=1$. Two non-zero coefficients, so the state is entangled, though less than maximally: the bases are the eigenvectors of $\rho^A$ and of $\rho^B$.
\end{solution}

\subsection{Purification}

The Schmidt decomposition shows that mixed states and entanglement are two sides of one coin. In the converse direction, every mixed state can be thought of as part of a pure state.

\begin{proposition}[Purification]\label{prop:purification}
  Let $\rho^A$ be a density operator on $\Hil_A$. Then there exists a pure state $\ket{AR}$ of a composite system $AR$, with $R$ a copy of $A$, such that $\Tr_R\ket{AR}\!\bra{AR}=\rho^A$.
\end{proposition}
\begin{proof}
  Write $\rho^A=\sum_ip_i\ket{i_A}\!\bra{i_A}$ in its eigenbasis and let $\{\ket{i_R}\}$ be an orthonormal basis of $R$. Define $\ket{AR}=\sum_i\sqrt{p_i}\ket{i_A}\ket{i_R}$. This is a unit vector, and by the computation in Corollary~\ref{cor:schmidt}(a), $\Tr_R\ket{AR}\!\bra{AR}=\sum_ip_i\ket{i_A}\!\bra{i_A}=\rho^A$.
\end{proof}

The state $\ket{AR}$ is a \term{purification} of $\rho^A$. It is not unique, but any two are related by a unitary acting on $R$ alone (Exercise~\ref{ex:purif-freedom}). The idea of purification is to treat the randomness of a mixed state as ignorance about a larger system, an \emph{environment} $R$ with which $A$ is entangled.
It lets us study mixed states using only state vectors, and it is the standard way to model noise in quantum computers.

\begin{example}\label{ex:purify}
  Purify the maximally mixed state $\tfrac12\Id$ and $\rho=\tfrac34\ket0\!\bra0+\tfrac14\ket1\!\bra1$.
\end{example}
\begin{solution}
  For $\tfrac12\Id=\tfrac12\ket0\!\bra0+\tfrac12\ket1\!\bra1$ the construction gives $\frac1{\sqrt2}(\ket{00}+\ket{11})$, a Bell state, as Example~\ref{ex:reduced-bell} suggests. For $\rho$, $\ket{AR}=\tfrac{\sqrt3}{2}\ket{00}+\tfrac12\ket{11}$: indeed $\Tr_R\ket{AR}\!\bra{AR}=\tfrac34\ket0\!\bra0+\tfrac14\ket1\!\bra1$.
\end{solution}

Purifications also tell how ensembles are generated: if $\{p_i,\ket{\psi_i}\}$ is an ensemble with density operator $\rho$, then $\sum_i\sqrt{p_i}\ket{\psi_i}\ket{i_R}$ is a purification, and measuring $R$ in the basis $\{\ket{i_R}\}$ yields $i$ with probability $p_i$, leaving $A$ in $\ket{\psi_i}$ (Exercise~\ref{ex:ensemble-purif}).
Different measurements of $R$ generate different ensembles of $A$ with the same $\rho^A$, which is the unitary freedom we met earlier.

\begin{exerciseset}
  \exercise{1} Find the Schmidt decompositions of $\frac1{\sqrt2}(\ket{00}+\ket{11})$, $\ket+\ket0$ and $\frac{1}{\sqrt2}(\ket{01}+\ket{10})$. What are the Schmidt numbers?

  \exercise{2} Find the Schmidt coefficients of $\frac{\ket{00}+\ket{01}+\ket{10}+\ket{11}}2$ and of $\frac1{\sqrt{5}}(\ket{00}+2\ket{11})$.

  \exercise{2} Compute the reduced density operator $\rho^A$ and its eigenvalues for $\frac1{\sqrt3}(\ket{00}+\ket{01}+\ket{10})$ (Example~\ref{ex:reduced-3term}), and find the Schmidt bases by diagonalising $\rho^A$.

  \exercise{2} Prove Corollary~\ref{cor:schmidt}: part (a), that the reduced operators share their spectrum, and that $\ket\psi$ is a product iff its Schmidt number is $1$. Show also that the Schmidt number is $\rank C$.

  \exercise{3} Let $\ket\psi,\ket\varphi$ be pure states of $AB$ with identical Schmidt coefficients. Show that there are unitaries $U$ on $A$ and $V$ on $B$ with $\ket\varphi=(U\tp V)\ket\psi$. Conclude that Schmidt coefficients are the invariants of the entanglement under local operations.

  \exercise{3} Extend the proof of Theorem~\ref{thm:schmidt} to the case $\dim\Hil_A\neq\dim\Hil_B$.

  \exercise{3} Show that a state of three systems $ABC$ need not have a decomposition $\ket\psi=\sum_i\lambda_i\ket{i_A}\ket{i_B}\ket{i_C}$ with orthonormal families. (Hint: for $\ket W=\frac1{\sqrt3}(\ket{001}+\ket{010}+\ket{100})$ compute $\rho^{AB}$; a state of the above form would have $\rho^{AB}=\sum_i\lambda_i^2\ket{i_Ai_B}\!\bra{i_Ai_B}$, whose eigenvectors are product states.)

  \exercise{2} Purify $\rho=\tfrac23\ket0\!\bra0+\tfrac13\ket1\!\bra1$, and the state $\rho=\tfrac12(\Id+\tfrac35\sz)$.

  \exercise{3} Let $\ket{AR_1},\ket{AR_2}$ be two purifications of $\rho^A$ to $AR$. Show that there is a unitary $U_R$ on $R$ with $\ket{AR_2}=(\Id\tp U_R)\ket{AR_1}$. (Hint: compare their Schmidt decompositions.)\label{ex:purif-freedom}

  \exercise{3} Let $\{p_i,\ket{\psi_i}\}$ be an ensemble with $\rho=\sum_ip_i\ket{\psi_i}\!\bra{\psi_i}$, and $R$ a system with orthonormal basis $\{\ket{i}\}$. (1) Show that $\sum_i\sqrt{p_i}\ket{\psi_i}\ket i$ is a purification of $\rho$. (2) Measuring $R$ in the basis $\{\ket i\}$ and getting outcome $i$, with what probability does this occur, and what is the corresponding state of $A$?\label{ex:ensemble-purif}
\end{exerciseset}

%======================================================================
\section{EPR and the Bell inequality}\label{sec:bell}
%======================================================================

\begin{objectives}
  You will explain the argument of Einstein, Podolsky and Rosen, derive the CHSH form of Bell's inequality from the assumption that
  physical properties exist before measurement and cannot be influenced at a distance, and show that quantum mechanics violates it.
\end{objectives}

\begin{recall}
  The singlet and the Bell states (Section~\ref{sec:composite}); expectation values of observables (Section~\ref{sec:measurement}); reduced density operators (Section~\ref{sec:density}); the groups $SU(2)$ and the representation $U\otimes U$ (Section~\ref{sec:grouprep}).
\end{recall}

% \begin{quote}
%   \emph{Anybody who is not shocked by quantum theory has not understood it.} --- Niels Bohr
% \end{quote}

We have built a mathematical framework and used it to compute. This last section asks what it \emph{means}. It is the point at which quantum mechanics stops being a convenient calculus and becomes a statement about the world that cannot be reproduced classically.

\subsection{The EPR argument}

In everyday life we assume that objects have properties independent of observation: the position of a tennis ball is there whether or not light bounces off it, and a measurement merely \emph{reveals} it.
Quantum mechanics seems to say otherwise. A qubit in the state $\ket0$ has no definite value of $X$: that value ``arises'' when $X$ is measured. Einstein found this unacceptable. In 1935, with Podolsky and Rosen~\cite{epr1935}, he argued that quantum mechanics must be \emph{incomplete}.

Their criterion for reality was the following: if we can predict with certainty the value of a physical quantity without disturbing the system, there is an \emph{element of reality} corresponding to that quantity, and any complete theory must contain it. The argument then uses an entangled state.
Take the singlet $\ket{\Psi^-}=\bell11=\frac1{\sqrt2}(\ket{01}-\ket{10})$, a state that is not a product (Example~\ref{ex:entangled-or-not}(d)).

\begin{proposition}[Perfect anti-correlation]\label{prop:singlet-anticorr}
  Let $\ket{\Psi^-}$ be the singlet and $\vec v$ any unit vector. If the observable $\vec v\cdot\vec\sigma$ is measured on both qubits, the two results are always opposite.
\end{proposition}
\begin{proof}
  \begin{steps}
    \item \emph{Local unitaries on the coefficient matrix.} For any unitary $U$, the coefficient matrix of $(U\tp U)\ket{\Psi^-}$ is $UCU^t$, with $C=\frac1{\sqrt2}\begin{pmatrix}0&1\\-1&0\end{pmatrix}$ (Exercise~\ref{ex:local-unitary-matrix}).
    \item \emph{A $2\times2$ identity.} $U\begin{pmatrix}0&1\\-1&0\end{pmatrix}U^t=(\det U)\begin{pmatrix}0&1\\-1&0\end{pmatrix}$, so
      \begin{equation}\label{eq:singlet-inv}
        (U\tp U)\ket{\Psi^-}=(\det U)\ket{\Psi^-}.
      \end{equation}
    \item \emph{Choose $U$ adapted to $\vec v$.} Let $\ket a,\ket b$ be the orthonormal eigenvectors of $\vec v\cdot\vec\sigma$ with eigenvalues $+1,-1$, and $U$ the unitary with $U\ket0=\ket a$, $U\ket1=\ket b$. Then $(U\tp U)\ket{\Psi^-}=\frac1{\sqrt2}(\ket{ab}-\ket{ba})$.
    \item \emph{Compare.} By \eqref{eq:singlet-inv} this equals $(\det U)\ket{\Psi^-}$, so up to a global phase $\ket{\Psi^-}=\frac1{\sqrt2}(\ket{ab}-\ket{ba})$.
    \item \emph{Read off the outcomes.} The outcomes $(+1,-1)$ and $(-1,+1)$ each have probability $\tfrac12$, and $(+1,+1),(-1,-1)$ have probability $0$. \qedhere
  \end{steps}
\end{proof}

Now put Alice and Bob far apart, each with one qubit of the singlet. Alice measures along $\vec v$ and gets $+1$. She can then predict with certainty that Bob, if he measures along the same $\vec v$, will get $-1$, without having touched his qubit. By EPR's criterion, the value of $\vec v\cdot\vec\sigma$ for Bob's qubit is an element of reality, and this holds for \emph{every} direction $\vec v$.
But quantum mechanics, as we have formulated it, contains no object representing the value of $\vec v\cdot\vec\sigma$ for all $\vec v$ at once: it only gives probabilities. Hence, EPR concluded, quantum mechanics is incomplete.

Most physicists were not convinced: imposing on Nature by fiat the properties she must obey is a peculiar way of studying her laws. And for thirty years the question seemed one of philosophy.

\subsection{Bell's inequality}

In 1964 Bell~\cite{bell1964} showed that EPR's world view has \emph{testable} consequences. We forget quantum mechanics for a moment and analyse a thought experiment with ``common sense''.

\begin{figure}[htb]
  \centering
  \begin{tikzpicture}[>=stealth]
    \node[draw=ink,thick,fill=ink!8,minimum width=1.6cm,minimum height=0.8cm,font=\small] (c) at (0,0) {Charlie};
    \node[draw=navy,thick,fill=navy!8,minimum width=2.6cm,minimum height=1.1cm,font=\small,align=center] (a) at (-4.6,0) {Alice\\$Q$ or $R$, $=\pm1$};
    \node[draw=navy,thick,fill=navy!8,minimum width=2.6cm,minimum height=1.1cm,font=\small,align=center] (b) at (4.6,0) {Bob\\$S$ or $T$, $=\pm1$};
    \draw[->,thick] (c) -- node[above,font=\footnotesize] {particle} (a);
    \draw[->,thick] (c) -- node[above,font=\footnotesize] {particle} (b);
    \draw[<->,figacc,dashed] (a.south) .. controls +(0,-1.2) and +(0,-1.2) .. node[below,font=\footnotesize] {no signal can connect the two choices} (b.south);
  \end{tikzpicture}
  \caption{The Bell experiment. Charlie prepares two particles and sends one to each of Alice and Bob, who are far apart. Each independently chooses one of two measurements, and each measurement has outcome $+1$ or $-1$.}
  \label{fig:bell-setup}
\end{figure}

Charlie prepares two particles, repeatably, and sends one to Alice and one to Bob (Figure~\ref{fig:bell-setup}). Alice randomly picks one of two measurements and Bob likewise: Alice measures a property $P_Q$ or $P_R$, Bob a property $P_S$ or $P_T$, each with outcome $\pm1$.
They do so at the same time, or rather in a causally disconnected way, so that by relativity Alice's measurement cannot influence Bob's and vice versa. We now make two assumptions that sound innocent.

\begin{itemize}
  \item \emph{Realism}: the properties have definite values $Q,R,S,T\in\{\pm1\}$ before and independently of the measurement, which merely reveals the one chosen.
  \item \emph{Locality}: Alice's choice and result do not influence Bob's result, and vice versa.
\end{itemize}
Together these two are called \term{local realism}. They are exactly the assumptions of the classical world view.

\begin{theorem}[CHSH inequality]\label{thm:chsh}
  Under local realism, for any preparation procedure,
  \[
    \mathbb E(QS)+\mathbb E(RS)+\mathbb E(RT)-\mathbb E(QT)\ \le\ 2 .
  \]
\end{theorem}
\begin{proof}
  \begin{steps}
    \item \emph{Factor the combination.} $QS+RS+RT-QT=(Q+R)S+(R-Q)T$.
    \item \emph{One term vanishes.} Since $Q,R=\pm1$, either $Q+R=0$ or $R-Q=0$, and the other bracket is $\pm2$. As $S,T=\pm1$, in every single run $QS+RS+RT-QT=\pm2$.
    \item \emph{Average over runs.} Let $p(q,r,s,t)$ be the probability that before the measurements $Q=q,R=r,S=s,T=t$; it may depend on Charlie's procedure and on noise. Then $\EE(QS+RS+RT-QT)=\sum_{qrst}p(q,r,s,t)(qs+rs+rt-qt)\le2\sum_{qrst}p(q,r,s,t)=2$.
    \item \emph{Split by linearity.} The left side equals $\EE(QS)+\EE(RS)+\EE(RT)-\EE(QT)$, and Alice and Bob can estimate each of these four averages from their data by pairing the runs in which they chose the corresponding measurements. \qedhere
  \end{steps}
\end{proof}

Notice what the theorem is: a statement about \emph{any} theory satisfying local realism, and not about quantum mechanics.

\begin{example}[Deterministic assignments]\label{ex:deterministic}
  Evaluate $QS+RS+RT-QT$ for the assignments $(Q,R,S,T)=(1,1,1,1)$, $(1,-1,1,-1)$ and $(-1,1,-1,-1)$, and find the largest value any local-realistic strategy can reach.
\end{example}
\begin{solution}
  \begin{steps}
    \item $(1,1,1,1)$: $1+1+1-1=2$.
    \item $(1,-1,1,-1)$: $QS=1$, $RS=-1$, $RT=1$, $QT=-1$, so $1-1+1+1=2$.
    \item $(-1,1,-1,-1)$: $QS=1$, $RS=-1$, $RT=-1$, $QT=1$, so $1-1-1-1=-2$.
    \item \emph{All $16$ assignments} give $\pm2$ (proof of Theorem~\ref{thm:chsh}), and a random mixture of them can only average to something in $[-2,2]$. So no classical strategy beats $2$.
  \end{steps}
\end{solution}

\subsection{Quantum mechanics violates the inequality}

Now put quantum mechanics back. Charlie prepares the singlet $\ket{\Psi^-}$ and sends the first qubit to Alice and the second to Bob. They measure
\[
  Q=Z_1,\quad R=X_1,\quad S=\frac{-Z_2-X_2}{\sqrt2},\quad T=\frac{Z_2-X_2}{\sqrt2}.
\]
The relevant expectation values come from a neat lemma.

\begin{lemma}\label{lem:singlet-corr}
  For real vectors $\vec a,\vec b$ and the singlet, $\ \bra{\Psi^-}(\vec a\cdot\vec\sigma)\tp(\vec b\cdot\vec\sigma)\ket{\Psi^-}=-\,\vec a\cdot\vec b$.
\end{lemma}
\begin{proof}
  \begin{steps}
    \item \emph{The singlet is invariant.} Put $U=e^{-\ii t\sigma_j}$ in \eqref{eq:singlet-inv}: $\det U=1$ because $\Tr\sigma_j=0$, so $(U\tp U)\ket{\Psi^-}=\ket{\Psi^-}$ for all $t$. (In the language of Section~\ref{sec:grouprep}: the singlet is fixed by the representation $U\tp U$ of $SU(2)$.)
    \item \emph{Differentiate at $t=0$.} $(\sigma_j\tp\Id+\Id\tp\sigma_j)\ket{\Psi^-}=0$, that is, $(\Id\tp\sigma_k)\ket{\Psi^-}=-(\sigma_k\tp\Id)\ket{\Psi^-}$.
    \item \emph{Move $\sigma_k$ across.} $\bra{\Psi^-}\sigma_j\tp\sigma_k\ket{\Psi^-}=-\bra{\Psi^-}\sigma_j\sigma_k\tp\Id\ket{\Psi^-}$.
    \item \emph{Reduce to the first qubit.} By Proposition~\ref{prop:reduced} and $\rho^A=\tfrac12\Id$ (Example~\ref{ex:reduced-bell}), this is $-\Tr(\sigma_j\sigma_k\,\rho^A)=-\tfrac12\Tr(\sigma_j\sigma_k)=-\delta_{jk}$.
    \item \emph{Linearity} gives the lemma for arbitrary $\vec a,\vec b$. \qedhere
  \end{steps}
\end{proof}

In the singlet the correlation of two spin measurements is minus the cosine of the angle between the axes. With $\vec z,\vec x$ the unit vectors of the axes, Alice measures along $\vec a_Q=\vec z$, $\vec a_R=\vec x$ and Bob along
$\vec b_S=-\frac{\vec z+\vec x}{\sqrt2}$, $\vec b_T=\frac{\vec z-\vec x}{\sqrt2}$ (Figure~\ref{fig:bell-axes}). By the lemma,
\[
  \langle QS\rangle=\tfrac1{\sqrt2},\quad \langle RS\rangle=\tfrac1{\sqrt2},\quad \langle RT\rangle=\tfrac1{\sqrt2},\quad \langle QT\rangle=-\tfrac1{\sqrt2},
\]
(for example $\langle QS\rangle=-\vec z\cdot\vec b_S=\frac1{\sqrt2}$ and $\langle QT\rangle=-\vec z\cdot\vec b_T=-\frac1{\sqrt2}$), and therefore
\[
  \langle QS\rangle+\langle RS\rangle+\langle RT\rangle-\langle QT\rangle=\tfrac4{\sqrt2}=2\sqrt2\approx2.83>2 .
\]

\begin{figure}[htb]
  \centering
  \begin{tikzpicture}[>=stealth,scale=1.5]
    \draw[->,ink!60] (-1.6,0) -- (1.6,0) node[right,font=\small] {$x$};
    \draw[->,ink!60] (0,-1.6) -- (0,1.6) node[above,font=\small] {$z$};
    \draw[very thick,navy,->] (0,0) -- (0,1.2) node[above right,font=\small] {Alice $Q$};
    \draw[very thick,navy,->] (0,0) -- (1.2,0) node[above right,font=\small] {Alice $R$};
    \draw[very thick,figacc,->] (0,0) -- (-0.85,-0.85) node[below left,font=\small] {Bob $S$};
    \draw[very thick,figacc,->] (0,0) -- (-0.85,0.85) node[above left,font=\small] {Bob $T$};
    \draw[ink!40] (0.5,0) arc (0:90:0.5);
  \end{tikzpicture}
  \caption{The four measurement axes in the $x$--$z$ plane. Bob's axes bisect the angles between Alice's, which maximises the violation of the CHSH inequality.}
  \label{fig:bell-axes}
\end{figure}

This contradicts Theorem~\ref{thm:chsh}. (One can show that $2\sqrt2$ is the largest value quantum mechanics allows, a result called Tsirelson's bound; Exercise~\ref{ex:tsirelson}.) Moreover the prediction has been tested: experiments with photons, starting with Aspect and collaborators~\cite{aspect1982}, and later ones
closing the experimental loopholes~\cite{hensen2015}, are overwhelmingly in favour of the quantum value. \emph{Nature does not obey the Bell inequality.}

\begin{example}\label{ex:chsh-numeric}
  Verify $\langle QS\rangle=\tfrac1{\sqrt2}$ directly, with the observables above and $\ket{\Psi^-}$.
\end{example}
\begin{solution}
  $QS=Z_1\cdot\frac{-Z_2-X_2}{\sqrt2}=-\tfrac1{\sqrt2}(Z\tp Z+Z\tp X)$. For the singlet $\langle Z\tp Z\rangle=-1$ (opposite results) and $\langle Z\tp X\rangle=0$ (Lemma~\ref{lem:singlet-corr} with $\vec a\perp\vec b$). Hence $\langle QS\rangle=-\tfrac1{\sqrt2}(-1+0)=\tfrac1{\sqrt2}$.
\end{solution}

\subsection{What it means}

Something in the derivation of Theorem~\ref{thm:chsh} must be wrong, and there were only two assumptions: realism and locality. Most physicists conclude that realism must go, but some prefer to give up locality; the vast literature
analyses the exact assumptions needed. The consensus is that Nature is not locally realistic. Two cautions are in order.

\begin{itemize}
  \item Quantum mechanics does \emph{not} allow faster-than-light signalling. Alice's choice of measurement does not change Bob's reduced state, which is $\tfrac12\Id$ whatever she does (Exercise~\ref{ex:no-signalling}), so Bob sees only random outcomes; the correlations appear only when the two lists of results are \emph{compared}, which requires ordinary communication.
  \item The violation is not an anomaly of one state. Every entangled pure state violates some Bell inequality, and for us the lesson is constructive: entanglement is a genuinely new resource, beyond anything classical.
\end{itemize}

\begin{lookahead}
  Bell's theorem has become a practical tool. In device-independent cryptography, two parties who observe a CHSH value above $2$ can certify, without trusting their devices, that their outcomes are genuinely random and unknown to any eavesdropper. Entropy bounds of the kind studied in Chapter~\ref{ch:info} (Section~\ref{sec:vonneumann}) are then used to quantify how much an eavesdropper can possibly learn.
\end{lookahead}

For quantum computing the lesson has been, historically, vague but fruitful: there is something profoundly different about entangled states, and a lot of progress has come from asking ``what would a little entanglement buy me in this problem?''. We have already seen one answer, superdense coding (Section~\ref{sec:superdense}), and teleportation (Section~\ref{sec:teleportation}) is another. Quantum cryptography and the speed-up of quantum algorithms are the subject of the courses that follow.

\begin{exerciseset}
  \exercise{1} Verify that the singlet $\ket{\Psi^-}$ is entangled, and that measuring $Z$ on both qubits always gives opposite results. Repeat with $X$ on both qubits.

  \exercise{2} Prove \eqref{eq:singlet-inv}: for $U\in M_2(\C)$ one has $U\begin{pmatrix}0&1\\-1&0\end{pmatrix}U^t=\det U\begin{pmatrix}0&1\\-1&0\end{pmatrix}$.

  \exercise{2} Verify Lemma~\ref{lem:singlet-corr} directly for $(\vec a,\vec b)=(\vec z,\vec z)$, $(\vec x,\vec x)$, $(\vec z,\vec x)$, by computing in the basis $\ket{ij}$.

  \exercise{2} Rotate all four measurement axes of Figure~\ref{fig:bell-axes} by the same angle $\alpha$ in the $x$--$z$ plane. Show with Lemma~\ref{lem:singlet-corr} that the CHSH value is still $2\sqrt2$, and explain why in terms of \eqref{eq:singlet-inv}. What happens if only Bob's axes are rotated?

  \exercise{2} Assuming local realism, show that $\mathbb E(QS)+\mathbb E(RS)+\mathbb E(RT)-\mathbb E(QT)\ge-2$, and that each of the four ``sign choices'' $(\pm QS\pm\dots)$ gives an inequality with the same bound.

  \exercise{3} (Tsirelson's bound.) Let $Q=\vec q\cdot\vec\sigma$, $R=\vec r\cdot\vec\sigma$, $S=\vec s\cdot\vec\sigma$, $T=\vec t\cdot\vec\sigma$ with real unit vectors. Show that $(Q\tp S+R\tp S+R\tp T-Q\tp T)^2=4\,\Id\tp\Id-\bigl[Q,R\bigr]\tp\bigl[S,T\bigr]$, and deduce $\langle QS+RS+RT-QT\rangle\le2\sqrt2$ in every state.\label{ex:tsirelson}

  \exercise{2} For the state $\ket{\Phi^+}=\bell00$, find directions $\vec a_Q,\vec a_R,\vec b_S,\vec b_T$ that also violate the CHSH inequality, and compute the maximal violation you can obtain.

  \exercise{2} Show that the state $\ket0\ket0$ does not violate the CHSH inequality for any choice of $\pm1$ observables of the above type, computing the largest value that $\langle QS+RS+RT-QT\rangle$ takes.

  \exercise{3} Show that Alice's choice of measurement cannot change Bob's reduced state in the singlet: if Alice measures with any projective measurement and does not tell Bob the result, Bob's state is still $\tfrac12\Id$. Explain why this prevents faster-than-light signalling.\label{ex:no-signalling}

  \exercise{3} A ``local hidden variable'' model assigns to each run a shared variable $\lambda$ with probability $\mu(\lambda)$ and outcomes $A(\vec a,\lambda),B(\vec b,\lambda)\in\{\pm1\}$. Show that Theorem~\ref{thm:chsh} applies to the correlations $E(\vec a,\vec b)=\int A(\vec a,\lambda)B(\vec b,\lambda)\mu(\lambda)\dd\lambda$, and that therefore no such model reproduces $E(\vec a,\vec b)=-\vec a\cdot\vec b$.
\end{exerciseset}

\section*{Further reading}

\begin{itemize}
  \item Nielsen and Chuang~\cite{nielsenchuang2010}, Chapter~2: the source that this chapter follows most closely in its second half. Its treatment of POVMs, the density operator and the Schmidt decomposition is more complete than ours.
  \item Sakurai and Napolitano~\cite{sakurai2020} and Griffiths and Schroeter~\cite{griffiths2018}: the postulates from the point of view of physics, with the continuous systems (position, momentum) that we have avoided.
  \item Mermin~\cite{mermin2007}: a short and unusually clear introduction to quantum computing, with a very good account of Bell's theorem.
  \item Preskill~\cite{preskill1998}: free lecture notes with a rigorous treatment of measurement, entanglement and mixed states.
  \item Einstein, Podolsky and Rosen~\cite{epr1935}, Bell~\cite{bell1964} and Clauser, Horne, Shimony and Holt~\cite{chsh1969}: the three papers behind Section~\ref{sec:bell}. Worth reading in the original; Bell's is only six pages.
\end{itemize}

